\section{Structure-sheaf direct images and uniform non-torsion}
\label{sec:analytic-picard}

This section reconstructs the analytic inputs used by the candidate's
Proposition~10.8.  The proof differs at one important point from the
candidate's route: extension across the cusp is obtained from the intrinsic
nearby-cycle isomorphism of \cref{thm:p0-integral-specialization}, not from
the cellular collapse.  The two cusp maps are now independently proved equal
in \cref{thm:collapse-equals-specialization}; the argument below nevertheless
retains its logically independent intrinsic formulation.

Put
\[
 e=d\zeta _1\wedge d\zeta _2
\]
on the marked torus family.  Retain the local coordinates and notation of
\cref{sec:candidate-geometry}:
\[
s_j\circ g_j=\zeta_js_j,\qquad
(m_1,m_2)=(3,4),\qquad t_j=s_j^{m_j},\qquad
f^{-1}(p_j)=m_jS_j.
\]

\subsection{The Hodge line and the finite normal characters}

\begin{lemma}[determinant characters and the Hodge line]
\label{lem:hodge-line-audit}
\Verified\ Let
\[
 \mathfrak K=f_*\omega_{X/B}.
\]
Then \(\mathfrak K\) is invertible and
\[
 \mathfrak K\simeq\mathcal O_B(1).
\tag{9.1}
\]
More precisely,
\[
 \det R_{g_1}=-\frac1\tau,\qquad
 \det R_{g_2}=\frac1\tau,\qquad
 \det R_{g_0}=1,
\tag{9.2}
\]
and at the fixed points
\[
 \det R_{g_j}(z_j)=\zeta_j^{a_j},
 \qquad (a_1,a_2)=(2,1).
\tag{9.3}
\]
The form \(e\) generates \(\mathfrak K\) at the cusp, while a unit times
\(s_j^{-a_j}e\) generates it at \(p_j\).
\end{lemma}

\begin{proof}
Equation (2.3d) in \cref{prop:period-input-audit} gives
\[
 R_g(z)\Pi(z)\lambda=\Pi(gz)A_g\lambda.
\]
Using
\[
 A_1\widehat w=\widehat u-\widehat w,\quad
 A_1\widehat\delta=\widehat\delta,\qquad
 A_2\widehat w=-\widehat u+\widehat\delta,\quad
 A_2\widehat\delta=\widehat\delta,
\]
and writing \(\tau^\sharp=\tau(g_jz)\),
\(\mu^\sharp=\mu(g_jz)\), direct multiplication gives
\[
 R_{g_1}=
 \begin{pmatrix}\tau^\sharp-1&0\\ \mu^\sharp&1\end{pmatrix},
 \qquad
 R_{g_2}=
 \begin{pmatrix}-\tau^\sharp&0\\1-\mu^\sharp&1\end{pmatrix}.
\]
The transformation laws
\(\tau\circ g_1=(\tau-1)/\tau\) and
\(\tau\circ g_2=-1/\tau\) prove (9.2) at \(g_1,g_2\); the cusp formulas
give \(R_{g_0}=I\).  Since \(\widehat\delta\) is fixed, one eigenvalue is
\(1\).  At \(\tau(z_1)=\rho=e^{\pi i/3}\), the other is
\(\rho-1=e^{2\pi i/3}=\zeta_1^2\); at \(\tau(z_2)=i\), it is
\(-i=\zeta_2\).  This proves (9.3).

We determine the extension over a finite filling without suppressing the
ramification.  On \(J'_j\), the nowhere-zero form \(ds_j\wedge e\)
transforms by
\[
 (\widetilde g_j^{\log})^*(ds_j\wedge e)
 =\chi_j(s_j)\,ds_j\wedge e,
 \qquad \chi_j=\zeta_j\det R_{g_j}.
\tag{9.4}
\]
Iteration \(m_j\) times shows
\(\prod_{r=0}^{m_j-1}\chi_j(\zeta_j^rs)=1\).
Write \(\chi_j(s)=\chi_j(0)e^{\psi(s)}\), with \(\psi(0)=0\).
The product identity says that the power series of \(\psi\) has no terms
whose degree is divisible by \(m_j\).  Hence
\[
 \varphi(s)-\varphi(\zeta_js)=\psi(s)
\]
has a holomorphic solution, obtained coefficientwise by division by
\(1-\zeta_j^n\).  Therefore
\[
 s^{k_j}e^{\varphi(s)}ds\wedge e,
 \qquad k_j=m_j-1-a_j,
\tag{9.5}
\]
is invariant: \(\chi_j(0)=\zeta_j^{1+a_j}\), so invariance is equivalent
to \(k_j+1+a_j\equiv0\pmod {m_j}\).  We now distinguish the descended
canonical section from the generator of its direct image.  Let
\[
 q_j:J'_j\longrightarrow N_j,\qquad
 \widetilde f_j:J'_j\longrightarrow\Delta_j,\qquad
 \varpi_j:\Delta_j\longrightarrow D_j
\]
denote the finite quotient, the torus-family projection, and the cyclic base
quotient, respectively, so that
\[
 f_j\circ q_j=\varpi_j\circ\widetilde f_j,
 \qquad t_j=s_j^{m_j}.
\]
Put \(\mathcal A_j=\widetilde f_j^{-1}(0)=\{s_j=0\}\) and
\(G_j=\langle\widetilde g_j^{\log}\rangle\).  The action is free by
\cref{prop:candidate-finite-fillings}; hence \(q_j\) is finite \'etale
and \(q_j^{-1}(S_j)=\mathcal A_j\).

The invariant form in (9.5), now denoted
\[
 \sigma_j=s_j^{k_j}e^{\varphi(s_j)}ds_j\wedge e,
 \qquad k_j=m_j-1-a_j,
\tag{9.5a}
\]
descends uniquely to
\(\overline\sigma_j\in H^0(N_j,K_{N_j})\) with
\(q_j^*\overline\sigma_j=\sigma_j\).  Since \(e^{\varphi(s_j)}\) is a
unit and \(ds_j\wedge e\) is nowhere vanishing,
\[
 \operatorname{div}_{J'_j}(\sigma_j)=k_j\mathcal A_j.
\]
Because \(q_j\) is \'etale,
\[
 q_j^*\operatorname{div}_{N_j}(\overline\sigma_j)
 =\operatorname{div}_{J'_j}(q_j^*\overline\sigma_j)
 =k_j\mathcal A_j=q_j^*(k_jS_j).
\]
Faithfulness of finite surjective pullback on Cartier divisors therefore gives
\[
 \operatorname{div}_{N_j}(\overline\sigma_j)=k_jS_j.
\tag{9.5b}
\]
Thus (9.5a) is a canonical section with that divisor, not a nowhere-vanishing
frame when \(k_j>0\); explicitly, \(k_1=0\) and \(k_2=2\).

For the direct-image generator, let
\[
 \vartheta_j=(ds_j\wedge e)\otimes(dt_j)^{-1}
\tag{9.5c}
\]
be the nowhere-vanishing frame of
\[
 q_j^*\omega_{N_j/D_j}
 =K_{J'_j}\otimes(f_j\circ q_j)^*K_{D_j}^{-1}.
\]
Here \((dt_j)^{-1}\) is the pulled-back frame of the inverse canonical line
of \(D_j\), not division by the differential of \(s_j^{m_j}\).  Put
\(\vartheta'_j=e^{\varphi(s_j)}\vartheta_j\).  The defining equation for
\(\varphi\), together with
\(\chi_j(s)=\chi_j(0)e^{\psi(s)}\) and
\(\chi_j(0)=\zeta_j^{1+a_j}\), gives
\[
 (\widetilde g_j^{\log})^*\vartheta'_j
 =\zeta_j^{1+a_j}\vartheta'_j.
\tag{9.5d}
\]
Every holomorphic coefficient of a relative top form is pulled back from
\(\Delta_j\), because the fibres of \(\widetilde f_j\) are connected compact
complex tori.  Hence
\[
 \widetilde f_{j*}q_j^*\omega_{N_j/D_j}
 =\mathcal O_{\Delta_j}\vartheta'_j.
\]
For \(F(s_j)\in\mathcal O_{\Delta_j}\), the section
\(F(s_j)\vartheta'_j\) is \(G_j\)-invariant exactly when
\[
 F(\zeta_js_j)\zeta_j^{1+a_j}=F(s_j).
\tag{9.5e}
\]
Writing \(F(s_j)=\sum_{n\ge0}c_ns_j^n\), this says that \(c_n=0\) unless
\[
 n+1+a_j\equiv0\pmod {m_j},
 \qquad\text{equivalently}\qquad n\equiv k_j\pmod {m_j}.
\]
Consequently,
\[
 \left((\varpi_j)_*\widetilde f_{j*}
 q_j^*\omega_{N_j/D_j}\right)^{G_j}
 =\mathcal O_{D_j}\cdot s_j^{k_j}\vartheta'_j,
\]
and finite descent gives the precise invariant-module identity
\[
 f_{j*}\omega_{N_j/D_j}
 =\mathcal O_{D_j}\cdot s_j^{k_j}\vartheta'_j.
\tag{9.5f}
\]
Finally, over \(\Delta_j^*\), the differential identity
\[
 dt_j=m_js_j^{m_j-1}ds_j
\]
identifies this generator with
\[
 s_j^{k_j}\vartheta'_j
 =\frac{e^{\varphi(s_j)}}{m_j}
 s_j^{k_j-(m_j-1)}e
 =\frac{e^{\varphi(s_j)}}{m_j}s_j^{-a_j}e.
\tag{9.5g}
\]
Thus a unit times \(s_j^{-a_j}e\) generates
\(f_{j*}\omega_{N_j/D_j}=(f_*\omega_{X/B})|_{D_j}\), whereas the
canonical section (9.5a) itself has divisor \(k_jS_j\).

At the cusp, in every smooth toric chart the unimodular monomial change of
coordinates turns
\[
 \Omega_0=\frac{dx_1}{x_1}\wedge
          \frac{dx_2}{x_2}\wedge dt_c
\]
into a unit times the ordinary coordinate volume form.  Thus it generates
the canonical line.  It is invariant under translations and under the
unimodular shears defining the quotient.  Under the punctured-overlap map,
\(\Omega_0=(2\pi i)^2dt_c\wedge e\), so \(e\) generates
\(\mathfrak K\) at \(p_0\).

It remains to compute the degree.  The modular-form calculation already
used in the proof of \cref{prop:period-input-audit} gives
\[
 \Theta=
 \left(\frac{E_4^2E_6^{1/2}}{\Delta_{\rm mod}}\right)\!\circ\tau.
\]
The square root exists because the zeros of \(E_6\circ\tau\) have even
order.  The weights \(8,3,12\) give the same automorphy factor as \(e\);
the sign at \(z_2\) is fixed by its linearizing coordinate.  The orders are
\[
 \operatorname{ord}_{z_1}\Theta=2=a_1,\qquad
 \operatorname{ord}_{z_2}\Theta=1=a_2,\qquad
 \Theta=t_c^{-1}\cdot(\text{unit})
\tag{9.6}
\]
at the cusp.  Consequently \(\mathsf e=\Theta^{-1}e\) is single-valued.
At \(p_j\) it is a unit times \(s_j^{-a_j}e\); at \(p_0\) it is a unit
times \(t_ce\).  Hence it is a global section of \(\mathfrak K\) with the
single simple zero \(p_0\).  This proves
\(\mathfrak K\simeq\mathcal O_B(p_0)\simeq\mathcal O_B(1)\).
Finally, \(\mathfrak K\) is torsion free of generic rank one on the smooth
curve \(B\), hence is invertible.
\end{proof}

\begin{lemma}[exact order of the finite-fibre normal line]
\label{lem:finite-normal-order-audit}
\Verified\ For \(j=1,2\),
\[
 \mathcal O_X(S_j)|_{S_j}=N_{S_j/X}
\]
has exact order \(m_j\) in \(\operatorname{Pic}(S_j)\).
\end{lemma}

\begin{proof}
On the finite cover of \cref{prop:candidate-finite-fillings}, the central
torus \(A_0^{(j)}=\{s_j=0\}\) maps to \(S_j\) as a connected principal
\(G_j\simeq\mathbb Z/m_j\)-cover.  Its normal line is trivial, and a
generator of \(G_j\) acts on the fibre by
\(s_j\mapsto\zeta_js_j\).  Thus \(N_{S_j/X}\) is associated to this
principal cover by the character \(\zeta_j\), and its \(m_j\)-th power is
trivial.

Suppose its \(k\)-th power were trivial for some \(0<k<m_j\).  Pulling a
nowhere-zero trivializing section back to \(A_0^{(j)}\) would give a
nowhere-zero holomorphic function \(h\) satisfying
\[
 h(gx)=\zeta_j^{k}h(x).
\]
Every holomorphic function on the compact connected torus \(A_0^{(j)}\)
is constant.  The displayed equation would force \(\zeta_j^k=1\), contrary
to \(0<k<m_j\).  Hence the order is exactly \(m_j\).
\end{proof}

% Complete finite-filling certificates, retaining the original marked periods.
% Host prerequisites: candidate_geometry.tex, exact_algebra.tex;
% normal-line comparison: analytic_picard.tex. No package or macro definitions.
\subsection{Exact product covers, normalization branches, and peripheral kernels of the finite fillings}
\label{ffcert:sec:finite-fillings}

We retain the actual period functions and finite fillings of
\cref{prop:period-input-audit,prop:monodromy-certificate,%
prop:candidate-finite-fillings}. The following explicit product maps
describe the original central tori and surfaces, not replacements of
their periods. We also identify the complete varying family as the
normalization of the actual ramified base change, compute the entire
peripheral inclusion kernel, and give a smooth trivialization of the
finite normal line. These calculations use no sphere identification,
global integral cohomology, or CDP theorem. The separate period and cusp
constructions retain their own proofs and are not inferred from these
finite-local calculations.

\subsubsection{The original lattice, periods, and their complex summands}

Use the ordered marked lattice and its dual basis
\[
 \Lambda=\mathbb Z\langle\widehat\gamma,\widehat u,
            \widehat w,\widehat\delta\rangle,
 \qquad \Lambda^*=\mathbb Z\langle\gamma,u,w,\delta\rangle.
 \tag{FF1}\label{ffcert:eq:marking}
\]
In these original column coordinates the matrices and translations are
\[
\begin{gathered}
 A_1=\begin{pmatrix}1&0&0&0\\6&0&1&0\\-6&-1&-1&0\\-2&1&0&1\end{pmatrix},
 \qquad
 A_2=\begin{pmatrix}1&0&0&0\\0&0&-1&0\\-6&1&0&0\\3&0&1&1\end{pmatrix},\\
 \varepsilon=(1,2,-4,0)^{\mathsf t},\quad
 \varepsilon'=(1,3,-3,0)^{\mathsf t},\quad
 v_1=\varepsilon,\quad v_2=-\varepsilon',\quad (m_1,m_2)=(3,4).
\end{gathered}
\tag{FF2}\label{ffcert:eq:matrices-translations}
\]
They satisfy $A_j^{m_j}=I$, $A_jv_j=v_j$, and
$\gamma A_j=\gamma$, with $\gamma(v_1)=1$ and $\gamma(v_2)=-1$.
Retain the entire period matrix
\[
 \Pi(z)=\begin{pmatrix}6\mu(z)&\tau(z)&1&0\\
                         \beta(z)&\mu(z)&0&1\end{pmatrix},
 \qquad
 D(z)=\operatorname{Im}\beta(z)
       -\frac{6(\operatorname{Im}\mu(z))^2}{\operatorname{Im}\tau(z)}<0.
 \tag{FF3}\label{ffcert:eq:periods}
\]
At the original elliptic centres let
\[
\begin{gathered}
 \rho=e^{\pi i/3},\quad \tau(z_1)=\rho,\quad \tau(z_2)=i,
 \quad \mu_1:=\mu(z_1)=\frac{2-\rho}{3},\quad
 \mu_2:=\mu(z_2)=\frac{1-i}{2},\\
 \beta_j=\beta(z_j),\qquad
 \kappa_1=\beta_1+2\mu_1,\qquad
 \kappa_2=\beta_2+3\mu_2.
\end{gathered}
\tag{FF4}\label{ffcert:eq:central-parameters}
\]
Neither $\beta_j$ nor either of its real components is adjusted.
Since $\operatorname{Im}\mu_1=-\operatorname{Im}\rho/3$ and
$\operatorname{Im}\mu_2=-1/2$, the exact inequalities are
\[
 \operatorname{Im}\kappa_1
 =\operatorname{Im}\beta_1-\tfrac23\operatorname{Im}\rho=D(z_1)<0,
 \qquad
 \operatorname{Im}\kappa_2
 =\operatorname{Im}\beta_2-\tfrac32=D(z_2)<0.
 \tag{FF5}\label{ffcert:eq:negative-orientation}
\]
Thus $(1,\kappa_j)$ is real-linearly independent with its original
negative imaginary orientation. No conjugation or positive-imaginary
replacement is made.

Put $\delta_*:=\widehat\delta$ and define
\[
\begin{aligned}
 e_1&=(0,1,-2,0)^{\mathsf t},&f_1&=(0,0,-3,1)^{\mathsf t},\\
 e_2&=(0,1,-1,0)^{\mathsf t},&f_2&=(0,0,-2,1)^{\mathsf t}.
\end{aligned}
\tag{FF6}\label{ffcert:eq:moving-generators}
\]

\begin{proposition}[The full fixed and moving integral lattices]
\label{ffcert:prop:integral-summands}
The exact fixed and image lattices are
\[
\begin{aligned}
 \Lambda^{A_1}&=\mathbb Z\varepsilon\oplus\mathbb Z\delta_*,&
 (A_1-I)\Lambda&=\mathbb Ze_1\oplus\mathbb Zf_1,\\
 \Lambda^{A_2}&=\mathbb Z\varepsilon'\oplus\mathbb Z\delta_*,&
 (A_2-I)\Lambda&=\mathbb Ze_2\oplus\mathbb Zf_2.
\end{aligned}
\tag{FF7}\label{ffcert:eq:integral-summands}
\]
The image lattices are saturated, with complete annihilators
\[
\begin{gathered}
 \operatorname{Ann}((A_1-I)\Lambda)
 =\mathbb Z\gamma\oplus\mathbb Z\psi_1,
 \quad \psi_1=2u+w+3\delta,\\
 \operatorname{Ann}((A_2-I)\Lambda)
 =\mathbb Z\gamma\oplus\mathbb Z\psi_2,
 \quad \psi_2=u+w+2\delta,\\
 \det[\varepsilon,\delta_*,e_1,f_1]=-3,
 \qquad \det[\varepsilon',\delta_*,e_2,f_2]=-2.
\end{gathered}
\tag{FF8}\label{ffcert:eq:annihilators-indices}
\]
In particular the sum of the fixed and moving lattices is not being
asserted to be all of $\Lambda$.
\end{proposition}

\begin{proof}
Solving $(A_1-I)(a,b,c,d)^{\mathsf t}=0$ gives $b=2a,c=-4a$;
the second equation gives $b=3a,c=-3a$. In both cases $a,d$ are arbitrary
integers. The columns of the two image matrices are, in order,
\[
 \begin{aligned}
 A_1-I &: \quad 6e_1-2f_1,\ -e_1+f_1,\ e_1,\ 0,\\
 A_2-I &: \quad 3f_2,\ -e_2,\ -e_2+f_2,\ 0.
 \end{aligned}
 \tag{FF9}\label{ffcert:eq:image-columns}
\]
These columns generate precisely the two stated image lattices. The first
is exactly the set of integer columns satisfying $a=0$ and
$2b+c+3d=0$: such a column is $be_1+df_1$. The second is exactly
the set satisfying $a=0$ and $b+c+2d=0$, with the analogous expression
$be_2+df_2$. They are therefore saturated. Their annihilators are the
displayed full integer spans: in the first case an annihilating row has
coefficients $(a,2c,c,3c)$, and in the second $(a,c,c,2c)$.
Expansion of the two four-by-four determinants gives $-3$ and $-2$.
\end{proof}

The complex images under the unchanged periods are
\[
\begin{aligned}
 \Pi(z_1)\varepsilon&=(0,\kappa_1),&
 \Pi(z_2)\varepsilon'&=(0,\kappa_2),&
 \Pi(z_j)\delta_*&=(0,1),\\
 \Pi(z_1)e_1&=(\rho-2,-(\rho-2)/3),&
 \Pi(z_1)f_1&=(-3,1),\\
 \Pi(z_2)e_2&=(i-1,-(i-1)/2),&
 \Pi(z_2)f_2&=(-2,1).
\end{aligned}
\tag{FF10}\label{ffcert:eq:complex-images}
\]
For example the first fixed coordinate vanishes because
$6(2-\rho)/3+2\rho-4=0$ or $6(1-i)/2+3i-3=0$, respectively.
The fixed real plane consequently maps onto the complex line
$\{(0,\xi):\xi\in\mathbb C\}$. The moving real planes map onto
$\{(\eta,-\eta/3)\}$ and $\{(\eta,-\eta/2)\}$, respectively,
because their displayed period pairs are real-linearly independent.
These are complementary complex lines in the original $\mathbb C^2$.

\subsubsection{The exact elliptic-product covers of the central tori and surfaces}

Define the one-dimensional complex tori
\[
\begin{gathered}
 E_j=\mathbb C/(\mathbb Z+\kappa_j\mathbb Z),\\
 F_1=\mathbb C/(\mathbb Z(\rho-2)+\mathbb Z(-3)),\qquad
 F_2=\mathbb C/(\mathbb Z(i-1)+\mathbb Z(-2)).
\end{gathered}
\tag{FF11}\label{ffcert:eq:elliptic-factors}
\]
Each period pair is real-linearly independent. Translation by its rank-two
lattice is free and properly discontinuous; the image of a closed real
parallelotope covers the quotient, proving compactness. Thus these are
smooth elliptic curves with the class of zero distinguished.

\begin{theorem}[Product maps onto the original finite central surfaces]
\label{ffcert:thm:product-covers}
The complex-linear maps
\[
 L_1(\xi,\eta)=(\eta,\xi-\eta/3),\qquad
 L_2(\xi,\eta)=(\eta,\xi-\eta/2)
 \tag{FF12}\label{ffcert:eq:addition-maps}
\]
induce surjective local biholomorphic group homomorphisms
\[
 \ell_j:E_j\times F_j\longrightarrow
 A_0^{(j)}:=\mathbb C^2/\Pi(z_j)\Lambda,
 \quad \ker\ell_1=\langle k_1\rangle,
 \quad \ker\ell_2=\langle k_2\rangle,
 \tag{FF13}\label{ffcert:eq:torus-covers}
\]
where
\[
 k_1=(2/3,\rho),\quad \operatorname{ord}(k_1)=3,
 \qquad k_2=(1/2,i),\quad \operatorname{ord}(k_2)=2.
 \tag{FF14}\label{ffcert:eq:kernel-points}
\]
The following automorphisms lift the original affine generators:
\[
 g^{\mathrm{prod}}_1(\xi,\eta)
   =(\xi+\kappa_1/3,(\rho-1)\eta),\qquad
 g^{\mathrm{prod}}_2(\xi,\eta)
   =(\xi-\kappa_2/4,-i\eta).
 \tag{FF15}\label{ffcert:eq:product-generators}
\]
Writing $k_j$ also for translation by that point, their joint groups
are $\langle k_1,g^{\mathrm{prod}}_1\rangle\cong C_3\times C_3$
and $\langle k_2,g^{\mathrm{prod}}_2\rangle\cong C_2\times C_4$.
They act freely, and the original reduced central surfaces are exactly
\[
 S_1=(E_1\times F_1)/\langle k_1,g^{\mathrm{prod}}_1\rangle,
 \qquad
 S_2=(E_2\times F_2)/\langle k_2,g^{\mathrm{prod}}_2\rangle.
 \tag{FF16}\label{ffcert:eq:surface-product-covers}
\]
The finite \'etale product-cover degrees are $9$ and $8$.
\end{theorem}

\begin{proof}
Each $L_j$ has determinant $-1$ in the displayed product and original
$(\zeta_1,\zeta_2)$ coordinates. Equation~\eqref{ffcert:eq:complex-images}
shows that the product period lattice maps exactly to the sum of the
fixed and moving sublattices inside $\Pi(z_j)\Lambda$.
Thus $L_j$ induces the stated surjective local biholomorphism, and its
kernel is identified with
$\Lambda/(\Lambda^{A_j}+(A_j-I)\Lambda)$, of order $3$ or $2$ by
\eqref{ffcert:eq:annihilators-indices}. Explicitly,
\[
 \widehat u=\tfrac23\delta_*+e_1-\tfrac23f_1,
 \qquad
 \widehat u=\tfrac12\delta_*+e_2-\tfrac12f_2.
 \tag{FF17}\label{ffcert:eq:kernel-lifts}
\]
Consequently $L_j(k_j)=\Pi(z_j)\widehat u$. The first kernel point
has order exactly three, since its fixed coordinate does, and
$3\rho=3(\rho-2)-2(-3)$ is a moving period. The second has order
exactly two, since its fixed coordinate does, and
$2i=2(i-1)-(-2)$ is a moving period. These orders equal the full kernel
cardinalities, proving the complete kernel assertions.

At the centres the original covariance matrices are
\[
 R_1=\begin{pmatrix}\rho-1&0\\(2-\rho)/3&1\end{pmatrix},\qquad
 R_2=\begin{pmatrix}-i&0\\(1+i)/2&1\end{pmatrix}.
 \tag{FF18}\label{ffcert:eq:central-covariance}
\]
They act by $1$ on the fixed lines and by
$\rho-1=e^{2\pi i/3}$ or $-i$ on the corresponding moving line.
Substitution gives the full equality of affine maps on $\mathbb C^2$,
\[
 L_jg_j^{\mathrm{prod}}
   =R_jL_j+\Pi(z_j)v_j/m_j.
 \tag{FF19}\label{ffcert:eq:product-intertwining}
\]
In particular the negative translation $-\kappa_2/4$ is forced by
$v_2=-\varepsilon'$. It has not been replaced by its opposite.

The moving rotations preserve their exact period lattices. For example,
using $\rho^2=\rho-1$, their images of the two period generators are
\[
\begin{aligned}
 (\rho-1)(\rho-2)&=-2(\rho-2)+(-3),&
 (\rho-1)(-3)&=-3(\rho-2)+(-3),\\
 (-i)(i-1)&=(i-1)-(-2),&
 (-i)(-2)&=2(i-1)-(-2).
\end{aligned}
\tag{FF20}\label{ffcert:eq:moving-rotations}
\]
The first matrix in these period bases has columns $(-2,1)$ and
$(-3,1)$, determinant $1$; the second has columns $(1,-1)$ and
$(2,-1)$, determinant $1$.
Equivalently, invariance and invertibility follow by applying the finite
orders of the rotations to these integral period images.
The kernel classes are fixed by the moving rotations because
\[
 (\rho-1)\rho-\rho=-1-\rho=-(\rho-2)+(-3),\qquad
 (-i)i-i=1-i=-(i-1).
 \tag{FF21}\label{ffcert:eq:kernel-rotation}
\]
Hence the product automorphism commutes with translation by $k_j$ on
the product torus. It has exact order $m_j$: its moving rotation has
that order; its $m_j$-th fixed-coordinate translation is
$\kappa_1$ or $-\kappa_2$; and no smaller positive multiple of that
translation is a fixed-coordinate period, since $\kappa_j$ is nonreal.

To prove both the direct-product assertion and freeness of the entire
joint group, put $(d_1,d_2)=(3,2)$, $(c_1,c_2)=(2/3,1/2)$ and
$(\epsilon_1,\epsilon_2)=(1,-1)$. A fixed point for $k_j^a
(g_j^{\mathrm{prod}})^b$, with $0\le a<d_j$ and $0\le b<m_j$,
would require
\[
 ac_j+\epsilon_jb\kappa_j/m_j=n+q\kappa_j,
 \qquad n,q\in\mathbb Z.
 \tag{FF22}\label{ffcert:eq:joint-freeness}
\]
Imaginary parts give $\epsilon_jb/m_j=q$, hence $b=0$. The real
equality then gives $ac_j\in\mathbb Z$, hence $a=0$ in its stated
range. Every nonidentity pair is therefore fixed-point free; in
particular it is not the identity transformation. This proves the
claimed joint groups. Quotienting first by $k_j$ and then by the lifted
affine generator gives exactly the original $S_j$ of
\cref{prop:candidate-finite-fillings}, by
\eqref{ffcert:eq:product-intertwining}. The total covering degrees are
$d_jm_j=9,8$ and the actions are free, so the covers are \'etale.
\end{proof}

\subsubsection{The varying quotient and its full local analytic coordinates}

Retain the original linearizing coordinate $s_j$ on $\Delta_j$, with
its actual round-disc radius $r_j>0$, and the original phase
\[
 s_j(g_jz)=\zeta_js_j(z),\quad
 \zeta_j=e^{-2\pi i/m_j},\quad t_j=s_j^{m_j},
 \qquad |s_j|<r_j,\quad |t_j|<r_j^{m_j}.
 \tag{FF23}\label{ffcert:eq:disc-and-phase}
\]
The radius is not replaced by one. On the universal marked cover the
holomorphic affine map and the exact period covariance are
\[
\begin{gathered}
 \widehat G_j(z,\zeta)
  =\bigl(g_jz,R_{g_j}(z)\zeta+\Pi(g_jz)v_j/m_j\bigr),\\
 R_{g_j}(z)\Pi(z)=\Pi(g_jz)A_j,
 \qquad t_\lambda(z,\zeta)=(z,\zeta+\Pi(z)\lambda).
\end{gathered}
\tag{FF24}\label{ffcert:eq:varying-action}
\]
Let $G_j$ denote the induced map of the original varying torus family
$J'_j=(\Delta_j\times\mathbb C^2)/\Lambda$ and put
$N_j=J'_j/\langle G_j\rangle$, with quotient map $q_j$.
The real-analytic marking $\Xi(z,x)=(z,\Pi(z)x)$ converts its powers
to the exact formulas
\[
 \Xi^{-1}\widehat G_j^k\Xi(z,x)
   =(g_j^kz,A_j^kx+kv_j/m_j),\qquad
 \widehat G_j^{m_j}=t_{v_j}.
 \tag{FF25}\label{ffcert:eq:varying-powers}
\]

\begin{proposition}[Smooth quotient charts and the exact multiplicities]
\label{ffcert:prop:quotient-charts}
The map $q_j:J'_j\to N_j$ is a finite \'etale covering of degree
$m_j$. The quotient is a Hausdorff smooth complex threefold, and its
proper map $f_j:N_j\to D_j$ has exact scheme fibre $m_jS_j$ at the
centre. In its local quotient coordinates the transverse map is
literally $t_j=s_j^{m_j}$.
\end{proposition}

\begin{proof}
At $s_j=0$ the freeness is proved in
\cref{ffcert:thm:product-covers}. Away from zero, a power with
$0<k<m_j$ has base coordinate multiplied by $\zeta_j^k\ne1$, so
has no fixed point. Equation~\eqref{ffcert:eq:varying-powers} shows
that the $m_j$-th power is the original lattice translation, hence is
the identity on $J'_j$. This gives a free action of exact order $m_j$.

For $y\in J'_j$, its finitely many distinct orbit points can be
separated by pairwise disjoint coordinate neighbourhoods. Intersecting
their inverse translates and shrinking gives a coordinate neighbourhood
$U$ of $y$ with the $G_j^aU$ pairwise disjoint. Then
$q_j|_U:U\to q_j(U)$ is a biholomorphism. Two distinct finite orbits
have disjoint saturated neighbourhoods by the same finite separation
argument. These charts prove the Hausdorff and smooth assertions, and
the quotient is a finite \'etale covering. They give directly the
trivial-stabilizer case of the quotient analytic structure in
\cite[conditions (I)--(II), Theorem~1]{Cartan1954Quotient}; smoothness
is not inferred from a theorem about arbitrary finite quotients.

At a point of $S_j$, use the quotient chart to retain $s_j$ as a
transverse coordinate, together with two fibre coordinates. The full
map is $t_j=s_j^{m_j}$. Thus $S_j$ is a smooth reduced Cartier
divisor, and its scheme multiplicity is exactly $m_j$.

For compact $C\subset\Delta_j$, the map induced by $\Xi$ from
$C\times[0,1]^4$ onto the inverse image of $C$ in $J'_j$ is
continuous and surjective, so that inverse image is compact. The actual
power map between the discs of radii $r_j$ and $r_j^{m_j}$ is proper:
the inverse image of a compact target set is closed, bounded, and stays
strictly inside the radius-$r_j$ disc. For compact $K\subset D_j$,
the inverse image in $J'_j$ of all roots of its $t_j$-values is
therefore compact. Its finite quotient image is exactly $f_j^{-1}(K)$,
proving properness of the full restriction.
\end{proof}

\subsubsection{Every branch in the normalization of the actual base change}

Let $\Delta_j^{\mathrm{base}}$ be the same radius-$r_j$ disc with
coordinate $u$, mapped to $D_j$ by $t_j=u^{m_j}$, and form
\[
 P_j=N_j\times_{D_j}\Delta_j^{\mathrm{base}},\qquad
 \nu_j:J'_j\longrightarrow P_j,\quad
 \nu_j(y)=(q_j(y),s_j(y)).
 \tag{FF26}\label{ffcert:eq:normalization-map}
\]

\begin{theorem}[The original varying family is the normalization]
\label{ffcert:thm:normalization}
The map $\nu_j$ is exactly the normalization of $P_j$. It is a
biholomorphism away from $u=0$; at the centre it has all $m_j$ local
branches prescribed by the original cyclic action. Its central fibre
is the connected original torus $A_0^{(j)}$, and
\[
 \nu_j(G_jy)=(q_j(y),\zeta_js_j(y)).
 \tag{FF27}\label{ffcert:eq:normalization-deck}
\]
Thus the original negative-phase action lifts $u\mapsto\zeta_ju$,
permutes the local branches cyclically, and has quotient $N_j$.
\end{theorem}

\begin{proof}
The definition uses no logarithm choice. If $u\ne0$, the $s_j$-values
in a covering orbit are exactly the $m_j$ distinct roots of its
$t_j$-value, so there is exactly one representative with value $u$.
If $u=0$, all $m_j$ covering points have $s_j=0$. This proves
surjectivity and describes all fibres. A compact subset of $P_j$
projects to a compact subset of $N_j$, whose inverse image under
$q_j$ is compact. Its inverse image under $\nu_j$ is a closed subset
of that compact set, hence is compact. Thus $\nu_j$ is proper with
finite fibres. Away from the centre the unique orbit representative
and quotient charts give the holomorphic inverse.

More precisely, around $n\in S_j$ choose a chart $q_j(U)$ as in
\cref{ffcert:prop:quotient-charts}, with inherited transverse coordinate
$s$ and fibre coordinates $z_1,z_2$. The fibre product is locally
\[
 s^{m_j}-u^{m_j}
   =\prod_{a=0}^{m_j-1}(s-e^{2\pi ia/m_j}u)=0.
 \tag{FF28}\label{ffcert:eq:all-branches}
\]
The factors are distinct smooth prime hypersurface factors in the
regular convergent-power-series ring $\mathbb C\{z_1,z_2,s,u\}$.
Their product ideal is radical: in this factorial ring, divisibility
of a power by every distinct prime factor implies divisibility of the
element by their product. Thus the fibre-product germ is reduced.
Its total ring of fractions is the product of the branch fraction
fields. Its integral closure is the product of the regular branch
rings $\mathbb C\{z_1,z_2,u\}$: adjoining the branch idempotents is
integral, since each satisfies $e^2-e=0$, and each component is
integrally closed. This proves the complete local normalization,
including all central branches.

On $G_j^aU$ the second component of $\nu_j$ is $u=\zeta_j^a s$,
or equivalently
\[
 s=\zeta_j^{-a}u=e^{2\pi ia/m_j}u.
 \tag{FF29}\label{ffcert:eq:branch-identification}
\]
Each of the $m_j$ disjoint covering charts therefore maps
isomorphically onto exactly one branch of
\eqref{ffcert:eq:all-branches}; these exhaust the branches. The
local identifications glue by the specified map $\nu_j$, giving the
global normalization. The source is the original varying family,
not an isotrivial product. Restricting its coordinate to zero gives
its connected torus $A_0^{(j)}$. Equation~\eqref{ffcert:eq:normalization-deck}
follows directly from $s_j(G_jy)=\zeta_js_j(y)$ and $q_jG_j=q_j$.
\end{proof}

\subsubsection{Universal covers and the full peripheral inclusion kernel}

\begin{theorem}[Exact finite-filling deck groups and attachment kernel]
\label{ffcert:thm:peripheral-kernel}
The universal cover $\Delta_j\times\mathbb C^2\to N_j$ has the
following complete deck-group presentation:
\[
 \Gamma_j=\left\langle\Lambda,\widehat G_j\ \middle|\
 \begin{array}{l}
 \Lambda\text{ abelian},\\
 \widehat G_jt_\lambda\widehat G_j^{-1}=t_{A_j\lambda},\\
 \widehat G_j^{m_j}=t_{v_j}
 \end{array}\right\rangle.
 \tag{FF30}\label{ffcert:eq:full-deck-group}
\]
This is also the universal-cover group of $S_j$. On the punctured
filling the complete group is $\Lambda\rtimes_{A_j}\mathbb Z$,
with the generator specified below, and the entire inclusion kernel is
\[
 \ker\{\pi_1(N_j\setminus S_j)\to\pi_1(N_j)\}
   =\langle\widetilde G_j^{m_j}t_{-v_j}\rangle
   \cong\mathbb Z.
 \tag{FF31}\label{ffcert:eq:whole-kernel}
\]
It is a central subgroup, not merely a necessary power relation.
\end{theorem}

\begin{proof}
In the real marking the original universal-cover transformations are
\[
 t_\lambda(s,x)=(s,x+\lambda),\qquad
 \widehat G_j(s,x)=(\zeta_js,A_jx+v_j/m_j).
 \tag{FF32}\label{ffcert:eq:cover-real-action}
\]
They satisfy the two relations in
\eqref{ffcert:eq:full-deck-group} by $A_jv_j=v_j$.
Moving translations to the left and using the power relation expresses
every word as $t_\lambda\widehat G_j^k$ with $\lambda\in\Lambda$
and $0\le k<m_j$. These forms are unique: the base action determines
$k$, and the remaining translation determines $\lambda$. Therefore
the abstract presented group acts faithfully and no additional relation
is implicit.

For $0<k<m_j$, a fixed-point equation would imply, on applying the
original integral covector $\gamma$ to the real fibre coordinate,
\[
 \gamma(\lambda)+k\gamma(v_j)/m_j=0.
 \tag{FF33}\label{ffcert:eq:cover-freeness}
\]
This is impossible because $\gamma(\lambda)$ is integral and
$\gamma(v_j)=1$ or $-1$. A nonzero translation with $k=0$ has no
fixed point either. The lattice action is properly discontinuous:
on compact base sets the real inverses $\Pi(z)^{-1}$ are uniformly
bounded, leaving finitely many lattice translations that can meet two
fixed compact sets. The full group is a finite union of lattice cosets,
so the same property holds for it. Thus the simply connected
$\Delta_j\times\mathbb C^2$ is the universal cover. Its restriction
to $s=0$ is the simply connected $\mathbb C^2$ covering $S_j$ with
the same action and group.

On the punctured disc retain exactly
\[
 w=\frac{\log s_j}{2\pi i},\qquad s_j=e^{2\pi iw},\qquad
 \operatorname{Im}w> -\frac{\log r_j}{2\pi},
 \qquad
 \widetilde G_j(w,x)=(w-1/m_j,A_jx+v_j/m_j).
 \tag{FF34}\label{ffcert:eq:log-cover-action}
\]
The displayed half-plane times $\mathbb C^2$, with its pulled-back
holomorphic coordinates, is simply connected. Along with the lattice
translations, $\widetilde G_j$ generates its entire punctured deck
group: the ordinary logarithm deck map is already
\[
 Q(w,x)=(w+1,x)=t_{v_j}\widetilde G_j^{-m_j}(w,x).
 \tag{FF35}\label{ffcert:eq:ordinary-log-deck}
\]
There is no finite-power relation for $\widetilde G_j$, since its
$k$-th power has $w$-displacement $-k/m_j$. Thus the exact normal
forms are $t_\lambda\widetilde G_j^k$ for unrestricted
$k\in\mathbb Z$, and the group is precisely the stated semidirect
product with action $A_j$.

The inclusion sends $t_\lambda$ to $t_\lambda$ and
$\widetilde G_j$ to $\widehat G_j$. Write $k=qm_j+r$ with
$0\le r<m_j$. Its image has normal form
$t_{\lambda+qv_j}\widehat G_j^r$, which is the identity exactly
when $r=0$ and $\lambda=-qv_j$. This proves
\eqref{ffcert:eq:whole-kernel} in both directions. Its generator has
$w$-displacement $-1$, so has infinite order. It commutes with all
translations because $A_j^{m_j}=I$, and with
$\widetilde G_j$ because $A_jv_j=v_j$, proving centrality.
\end{proof}

\subsubsection{The original punctured overlap and integral cohomology pullback}

On the affine punctured model choose a local logarithm and put
\[
 \alpha(z)=\frac{\log s_j(z)}{2\pi i},\qquad
 h_j[z,\zeta]=[z,\zeta+\alpha(z)\Pi(z)v_j].
 \tag{FF36}\label{ffcert:eq:overlap-map}
\]
A branch change adds $n\Pi(z)v_j$, an original lattice period,
so the map is well-defined and holomorphic on the punctured torus
quotient; subtraction of the same term gives its inverse. Under the
clockwise step $s_j\mapsto\zeta_js_j$, the compatible branch is
$\alpha(g_jz)=\alpha(z)-1/m_j$. The covariance
\eqref{ffcert:eq:varying-action} and $A_jv_j=v_j$ therefore give
the exact conjugacy to the original linear lift
$\widehat G_j^0(z,\zeta)=(g_jz,R_{g_j}(z)\zeta)$:
\[
 h_j\widehat G_jh_j^{-1}=\widehat G_j^0.
 \tag{FF37}\label{ffcert:eq:overlap-conjugacy}
\]
Indeed the two translations in $h_j\widehat G_j$ add to
$\Pi(g_jz)v_j(1/m_j+\alpha(g_jz))
=\alpha(z)\Pi(g_jz)v_j$, exactly the translation in
$\widehat G_j^0h_j$.

The corresponding real translation retains the full complex structure:
\[
\begin{gathered}
 J_z=\Pi(z)^{-1}i\Pi(z),\qquad
 a_j(z)=\operatorname{Re}\alpha(z)v_j
               +\operatorname{Im}\alpha(z)J_zv_j,\\
 J_{g_jz}A_j=A_jJ_z,\qquad
 a_j(g_jz)=A_ja_j(z)-v_j/m_j.
\end{gathered}
\tag{FF38}\label{ffcert:eq:real-overlap}
\]
This follows by applying $\Pi(z)^{-1}$ to
$\alpha(z)\Pi(z)v_j$ and then using covariance; complex scalar
multiplication is not being imposed on the marked real lattice.
The inverse image of the middle zero section is consequently
$-a_j(z)$, whose endpoint after one clockwise step is exactly
$A_j(-a_j(z))+v_j/m_j$. Iterating $m_j$ such steps increases
this marked endpoint by precisely $v_j$, not by an unspecified
translated generator. Thus the overlap map realizes the based
peripheral convention in \eqref{ffcert:eq:whole-kernel}.

The map from the fibre over $z$ to its deck-translated fibre over
$g_jz$ has marked linear part $A_j$; its inverse has part
$A_j^{-1}$. The original contragredient convention is
\[
 (A_j^{-1})^{\mathsf t}=T_j,
 \qquad A_j=(T_j^{-1})^{\mathsf t}.
 \tag{FF39}\label{ffcert:eq:contragredient-direction}
\]
These directions keep a based loop distinct from its inverse deck
map. No matrix or generator is reversed to obtain the attachment relation.

\begin{corollary}[The exact integral first-cohomology covering image]
\label{ffcert:cor:covering-h1}
For the original covering $\pi_j:A_0^{(j)}\to S_j$,
\[
 H_1(S_j,\mathbb Z)\cong\mathbb Z^2,\qquad
 \pi_j^*H^1(S_j,\mathbb Z)
   =\mathbb Z(m_j\gamma)\oplus\mathbb Z\psi_j
        \subset H^1(A_0^{(j)},\mathbb Z).
 \tag{FF40}\label{ffcert:eq:integral-h1-image}
\]
The index $m_j$ in the invariant covector lattice is retained.
\end{corollary}

\begin{proof}
By \cref{ffcert:prop:integral-summands}, the map
\[
 \Lambda\longrightarrow\mathbb Z^2,\qquad
 \lambda\longmapsto(\gamma(\lambda),\psi_j(\lambda))
 \tag{FF41}\label{ffcert:eq:coinvariant-coordinates}
\]
has kernel exactly $(A_j-I)\Lambda$. It is onto:
$\widehat\gamma$ gives $(1,0)$, and $\widehat w$ for $j=1$,
or $\widehat u$ for $j=2$, gives $(0,1)$. Thus the full
coinvariant group is $\mathbb Z^2$. In these coordinates $v_j$
is $(1,0)$ or $(-1,0)$. Abelianizing
\eqref{ffcert:eq:full-deck-group} adjoins one generator with relation
$m_j\widehat G_j=v_j$, eliminating a primitive generator and
leaving $\mathbb Z^2$. This computes $H_1(S_j,\mathbb Z)$ without
global topology. A homomorphism from this group to $\mathbb Z$
restricts on $\Lambda$ to $a\gamma+b\psi_j$ precisely when
$m_j\mid a$, since it must assign the integer
$a\gamma(v_j)/m_j$ to $\widehat G_j$. This necessary and sufficient
condition proves the exact covering pullback.
\end{proof}

\subsubsection{The finite-to-global analytic gluing interface}

The global construction of \cref{thm:candidate-global-geometry} glues
complete punctured restrictions of the original middle family and
the three local fillings. The finite identifications are precisely
\eqref{ffcert:eq:overlap-map}. The special discs are pairwise disjoint,
and a middle point over each punctured disc has exactly one
representative in its filling. Thus an equivalence class meets any
given summand at most once. Saturations of open subsets are open,
because the overlap identifications are biholomorphisms, so each summand
embeds openly and injectively.

Two points at one special base value have representatives in its same
filling; two points at one ordinary base value have representatives in
the middle piece. They have disjoint neighbourhoods in that Hausdorff
open summand, and their images remain disjoint because no equivalence
class meets that summand twice. Points at distinct base values are
separated by the base. This proves the gluing's Hausdorff assertion
using the actual full restrictions, not an unrestricted gluing principle.

In particular $f^{-1}(D_j)$ is the entire corresponding filling:
every middle point over $D_j^*$ is identified with its unique filling
representative. Thus the properness in
\cref{ffcert:prop:quotient-charts} is properness of the full inverse
image. For a compact base set $K$, finitely many relatively compact
$V_a$ whose closures lie inside members $U_a$ of the original
four-piece base cover can be chosen to cover $K$. The compact sets
$K\cap\overline V_a$ have compact preimages under the corresponding
full proper restrictions, and these preimages cover exactly
$f^{-1}(K)$. This is the compact-cover argument for global properness.
The separate cusp and middle constructions supply their own full
proper restrictions; the finite-local proof neither replaces those
constructions nor implies smooth sphere recognition.

\subsubsection{The exact normal character and a smooth trivialization}

\begin{theorem}[Holomorphic normal order and the explicit smooth product]
\label{ffcert:thm:normal-trivialization}
Let $\mathcal N_j=N_{S_j/N_j}$, equivalently
$\mathcal O_X(S_j)|_{S_j}$ after the specified open embedding into
the original global analytic construction. This line has exact order
$m_j$ in the holomorphic Picard group, but has the following explicit
smooth trivialization and hence zero integral first Chern class.
\end{theorem}

\begin{proof}
Differentiating the actual transverse coordinate in
\eqref{ffcert:eq:disc-and-phase} identifies the normal line as the
quotient of the trivial line over $A_0^{(j)}$ by
\[
 (x,n)\longmapsto(A_jx+v_j/m_j,\zeta_j n),
 \qquad \zeta_j=e^{-2\pi i/m_j}.
 \tag{FF42}\label{ffcert:eq:normal-action}
\]
Here $x$ is the original real marking of the central complex torus;
the action on the torus and line is holomorphic. A section of its
$k$-th power therefore pulls back to a function satisfying
\[
 h(A_jx+v_j/m_j)=\zeta_j^k h(x).
 \tag{FF43}\label{ffcert:eq:normal-section-character}
\]
The positive exponent follows from the quotient itself: the representative
$(x,h(x))$ is identified at the image point with
$(A_jx+v_j/m_j,\zeta_j^kh(x))$, which must equal the representative
of the section there. The negative exponent would instead describe
the dual normal line under the declared action.

The $m_j$-th power has the trivial character, so the constant function
one descends to a nowhere-zero holomorphic section of that power.
If a power $0<k<m_j$ were holomorphically trivial, its trivializing
section would pull back to a nowhere-zero holomorphic function on the
compact connected torus $A_0^{(j)}$. Such a function is constant,
and \eqref{ffcert:eq:normal-section-character} would force
$\zeta_j^k=1$, impossible in this range. This proves the exact
holomorphic order in \cref{lem:finite-normal-order-audit}, with the
normal tangent character fully specified.

For the smooth trivialization put
\[
 \epsilon_j=\gamma(v_j)\in\{1,-1\},\qquad
 H_j(x)=\exp(-2\pi i\epsilon_j\gamma(x))
 \quad(x\in\Lambda\otimes\mathbb R).
 \tag{FF44}\label{ffcert:eq:smooth-normal-section}
\]
Integrality of $\gamma(\lambda)$ gives $H_j(x+\lambda)=H_j(x)$,
so this is a well-defined nowhere-zero smooth function on the actual
marked central torus. Moreover $\gamma A_j=\gamma$ and
$\epsilon_j^2=1$ give
\[
 \begin{aligned}
 H_j(A_jx+v_j/m_j)
 &=\exp(-2\pi i\epsilon_j\gamma(x))
       \exp(-2\pi i\epsilon_j^2/m_j)\\
 &=\zeta_jH_j(x).
 \end{aligned}
 \tag{FF45}\label{ffcert:eq:smooth-equivariance}
\]
It consequently descends to a nowhere-zero section of the actual
normal line. Explicitly the smooth fibrewise-linear trivialization and
its inverse are
\[
 \begin{aligned}
 \Theta_j:\mathcal N_j&\longrightarrow S_j\times\mathbb C,
 & [x,n]&\longmapsto([x],n/H_j(x)),\\
 \Theta_j^{-1}:S_j\times\mathbb C&\longrightarrow\mathcal N_j,
 & ([x],z)&\longmapsto[x,zH_j(x)].
 \end{aligned}
 \tag{FF46}\label{ffcert:eq:normal-product-isomorphism}
\]
Replacing $x$ by a lattice translate leaves both formulas unchanged.
Replacing $(x,n)$ by
$(A_jx+v_j/m_j,\zeta_jn)$ also leaves $n/H_j(x)$ unchanged by
\eqref{ffcert:eq:smooth-equivariance}; the inverse has the same
equivariance. The two maps are therefore well-defined, smooth, and
inverse, retaining the original signs for each $j$. Thus
$c_1(\mathcal N_j)=0$ in integral cohomology while its holomorphic
order is exactly $m_j$. This conclusion uses only the original affine
action and integral covector, not the topology of a global $X$.
\end{proof}


\subsection{The cusp structure sheaf}

\begin{lemma}[structure-sheaf cohomology of the cusp fibre]
\label{lem:cusp-structure-sheaf-audit}
\Verified\ The reduced normal-crossings fibre \(W\) satisfies
\[
 h^q(W,\mathcal O_W)=(1,2,1)\qquad(q=0,1,2).
\tag{9.7}
\]
The natural map
\[
 H^1(W;\mathbb C)\longrightarrow H^1(W,\mathcal O_W)
\tag{9.8}
\]
is an isomorphism.  In degree two the corresponding map is surjective and
is injective on the kernel of pullback to the normalization
\(\nu:\mathrm{dP}_6\to W\).  Finally,
\[
 \bigwedge^2H^1(W,\mathcal O_W)
 \longrightarrow H^2(W,\mathcal O_W)
\tag{9.9}
\]
is an isomorphism.
\end{lemma}

\begin{proof}
Let \(D_1,D_2,D_3\simeq\mathbb P^1\) be the double curves and \(P,Q\) the
two triple points, with the branch order fixed in the proof of
\cref{thm:p0-integral-specialization}.  The normalization complex is
\[
0\longrightarrow\mathcal O_W\longrightarrow
 \nu_*\mathcal O_{\mathrm{dP}_6}\xrightarrow{\delta_0}
 \bigoplus_{k=1}^3\mathcal O_{D_k}\xrightarrow{\delta_1}
 \mathbb C_P\oplus\mathbb C_Q\longrightarrow0,
\tag{9.10}
\]
where \(\delta_0\) is the difference of the two branch restrictions and
\[
 \delta_1(g_1,g_2,g_3)
 =(g_1(P)-g_2(P)+g_3(P),
   g_1(Q)-g_2(Q)+g_3(Q)).
\]
For completeness, exactness is stalkwise.  At a generic double point it is
the fibre-product sequence
\[
0\to\mathbb C\{x,y,z\}/(xy)
 \to\mathbb C\{y,z\}\oplus\mathbb C\{x,z\}
 \to\mathbb C\{z\}\to0.
\]
At a triple point it is the augmented branch complex for
\(\mathbb C\{x,y,z\}/(xyz)\): three branch rings, three pairwise
intersection rings, and the common residue field, with alternating
differences.  The augmented two-simplex complex is exact; successively
lifting compatible power series proves exactness of the analytic rings.
Away from the singular locus the first arrow is an isomorphism.

The same formulas give an exact complex with \(\mathbb C\) in place of
\(\mathcal O\).  Since \(\mathrm{dP}_6\) is rational and each \(D_k\) is a
rational curve, the positive structure-sheaf cohomology of every term on
the right of (9.10) vanishes.  Hypercohomology therefore reduces to
\[
 \mathbb C\xrightarrow{0}\mathbb C^3
 \xrightarrow{d_1}\mathbb C^2,
 \qquad
 d_1(a_1,a_2,a_3)=(s,s),\quad s=a_1-a_2+a_3.
\tag{9.11}
\]
Its cohomology has dimensions \(1,2,1\), proving (9.7).

For the constant-sheaf complex, the row of degree zero is again (9.11),
the row of degree one vanishes, and the remaining degree-two row begins
with \(H^2(\mathrm{dP}_6;\mathbb C)\).  The comparison induced by
\(\mathbb C_W\hookrightarrow\mathcal O_W\) is the identity on (9.11) and
zero on the positive rows of the target.  It follows that (9.8) is an
isomorphism, while in degree two it maps the dual-complex summand
\(\ker\nu^*\) isomorphically onto \(H^2(W,\mathcal O_W)\).

Choose the basis \(\gamma_W,u_W\in H^1(W;\mathbb Z)\) whose intrinsic
specializations are \(\gamma,u\), respectively; it exists and is unique by
\cref{thm:p0-integral-specialization}.  Their cup product is nonzero,
because its specialization is
\(\gamma\wedge u\ne0\in H^2(F;\mathbb Z)\).  Its pullback to
\(\mathrm{dP}_6\) is zero since
\(H^1(\mathrm{dP}_6;\mathbb C)=0\).  The degree-two comparison therefore
sends it to a nonzero element of \(H^2(W,\mathcal O_W)\).  The images of
\(\gamma_W,u_W\) form a basis by (9.8), and both sides of (9.9) are
one-dimensional.  Hence (9.9) is an isomorphism.
\end{proof}

\subsection{Direct images, base change, and the Picard line}

\begin{theorem}[structure-sheaf direct images]
\label{thm:structure-direct-images-audit}
\Verified\ For the constructed map \(f:X\to B=\mathbb P^1\),
\[
 f_*\mathcal O_X=\mathcal O_B,\qquad
 R^1f_*\mathcal O_X\simeq
 \mathcal O_B\oplus\mathcal O_B(-1),\qquad
 R^2f_*\mathcal O_X\simeq\mathcal O_B(-1).
\tag{9.12}
\]
All three direct images are locally free and every base-change morphism
\[
 (R^qf_*\mathcal O_X)\otimes k(b)
 \longrightarrow H^q(X_b,\mathcal O_{X_b}),
 \qquad 0\le q\le2,
\tag{9.13}
\]
is an isomorphism, including at the two nonreduced fibres.  Moreover the
invariant class \(\gamma\) defines the trivial summand in
\(R^1f_*\mathcal O_X\) and is nowhere zero on \(B\).
\end{theorem}

\begin{proof}
The proper map has connected fibres and \(B\) is normal, so Stein
factorization gives \(f_*\mathcal O_X=\mathcal O_B\).

The map is flat.  A local parameter on \(B\) pulls back to a nonzero
element of the regular local domain \(\mathcal O_{X,x}\); if it vanished
identically near one point, analytic continuation on the connected
manifold would make \(f\) constant.  Thus \(\mathcal O_{X,x}\) is
torsion-free, hence flat, over the discrete valuation ring
\(\mathcal O_{B,f(x)}\).

Over \(B^\circ\), the fibres are two-dimensional complex tori and their
structure-sheaf cohomology dimensions are \(1,2,1\).  At the cusp the same
dimensions hold by \cref{lem:cusp-structure-sheaf-audit}.  Grauert's
theorem therefore gives local freeness of ranks \(1,2,1\) and base change
on a neighbourhood of \(B^\circ\cup\{p_0\}\); the precise analytic statement is
\cite[Chapter~III, \S4, Theorem~4.12]{BanicaStanasila1976}, originally due to
Grauert \cite{Grauert1960}.

At \(p_j\), use the finite cover \(J'_j\to N_j\).  If
\(\widetilde f:J'_j\to\Delta_j\) is the smooth torus family and
\(\pi(s)=s^{m_j}\), exactness of finite-group invariants in characteristic
zero gives
\[
 R^qf_*\mathcal O_{N_j}
 =\bigl(\pi_*R^q\widetilde f_*\mathcal O_{J'_j}\bigr)^{G_j}.
\tag{9.14}
\]
After averaging, a free eigenframe \(v\) may be chosen for each character
\(\nu\).  The invariant sections are exactly
\(s^rv\,\mathcal O_{D_j}\), where \(0\le r<m_j\) is determined by
\(\nu\zeta_j^r=1\).  Hence (9.14) is free of ranks \(1,2,1\).
Higher direct images vanish by the fibre-dimension bound, so these are all
the nonzero direct images.  Since all direct images are locally free, the analytic
cohomology-and-base-change theorem, applied by descending induction from
\(q=2\), proves (9.13) at the finite fibres as well.  In the notation of
\cite[Chapter~III, \S3, Theorem~3.4, Corollaries~3.5, 3.7, and 3.10]
  {BanicaStanasila1976}.  Corollary~3.10 says that, for each fixed \(q\),
  freeness of \(R^p f_*\mathcal O_X\) for every \(p\ge q\) is equivalent to
  cohomological flatness in all dimensions \(p\ge q\).  Applying this
  successively with \(q=2,1,0\) gives the displayed base-change maps in every
  relevant degree.  This is a local statement on the base and so applies
  equally to the nonreduced scheme-theoretic fibres.

Write \(G=R^1f_*\mathcal O_X\).  We construct its distinguished section
without using a cellular collapse.  On \(B^\circ\), the monodromy-invariant
integral class \(\gamma\) maps to a holomorphic section
\(\gamma^{0,1}\) of \(G\).  It is nowhere zero there: if its image in
\(H^{0,1}(F)\) vanished, then
\(\gamma\in H^{1,0}(F)\cap H^1(F;\mathbb R)=0\).  On a finite cover it is
a constant invariant class and remains nonzero at \(s_j=0\), so it
descends across \(p_j\).

At the cusp, let \(j:\Delta^*\hookrightarrow\Delta\).  The natural
specialization morphism
\[
 R^1f_*\mathbb C\longrightarrow
 j_*R^1f_*\mathbb C|_{\Delta^*}
\]
is an isomorphism: away from \(0\) this is tautological, while at \(0\) it
is exactly the isomorphism from \(H^1(W;\mathbb C)\) onto the
\(T_0\)-invariant subspace proved in
\cref{thm:p0-integral-specialization}.  Hence the invariant classes
\(\gamma,u\) extend as sections of \(R^1f_*\mathbb C\) on the whole disc.
The natural morphism induced by \(\mathbb C\hookrightarrow\mathcal O\),
\[
 R^1f_*\mathbb C\otimes\mathcal O_\Delta
 \longrightarrow G|_\Delta,
\]
sends them at the central point to the two classes in
\(H^1(W,\mathcal O_W)\).  These form a basis by
\cref{lem:cusp-structure-sheaf-audit}.  Thus
\(\gamma^{0,1}\) extends across the cusp and is nowhere zero everywhere.
In particular \(\mathcal O_B\gamma^{0,1}\subset G\) is a line subbundle.

Relative duality applies to the proper flat Gorenstein map \(f\): \(X\) is
smooth, the base is a smooth curve, and \(\omega_{X/B}\) is invertible.
Using \cref{lem:hodge-line-audit}, it gives
\[
 R^2f_*\mathcal O_X\simeq\mathfrak K^\vee
 \simeq\mathcal O_B(-1).
\tag{9.15}
\]
The relative theorem is the one of Ramis--Ruget--Verdier
\cite[pp.~261--283]{RamisRugetVerdier1971}; the absolute singular
dualizing-complex comparison is independently recorded in
\cite[Theorems~1--2 and Proposition~2]{RamisRuget1970}.

The relative cup product
\[
 \bigwedge^2G\longrightarrow R^2f_*\mathcal O_X
\tag{9.16}
\]
is an isomorphism.  It is so on every smooth torus fibre.  At the cusp it
is so by (9.13) and \cref{lem:cusp-structure-sheaf-audit}.  At \(p_j\),
choose on the cover an eigenframe \(v_1,v_2\), with
\(v_1=\gamma^{0,1}\) invariant.  Its invariant exponent is \(r_1=0\).
The action on \(v_1\wedge v_2\) is dual to the Hodge character in (9.3),
so its invariant generator is \(s_j^{a_j}(v_1\wedge v_2)\).
The invariant generator for the second line is consequently
\(s_j^{a_j}v_2\), because \(0<a_j<m_j\).  Wedge product sends
\(v_1\wedge(s_j^{a_j}v_2)\) to exactly the invariant generator of the
right side of (9.16), with no wrap modulo \(m_j\).

Consequently
\(\det G\simeq R^2f_*\mathcal O_X\simeq\mathcal O_B(-1)\).  The quotient
of \(G\) by its trivial subbundle is \(\mathcal O_B(-1)\), and the
extension splits because
\[
 \operatorname{Ext}^1(\mathcal O_B(-1),\mathcal O_B)
 =H^1(B,\mathcal O_B(1))=0.
\]
This proves (9.12) and identifies the trivial summand.
\end{proof}

\begin{corollary}[the Picard parameter and fibrewise restriction]
\label{cor:picard-line-audit}
\Verified\ One has
\[
 H^1(X,\mathcal O_X)=\mathbb C,\qquad
 \operatorname{Pic}(X)=\operatorname{Pic}^0(X)\simeq(\mathbb C,+).
\tag{9.17}
\]
For every scheme-theoretic fibre \(X_b\), restriction induces an injective
complex-linear map
\[
 \iota_b:H^1(X,\mathcal O_X)\longrightarrow
 H^1(X_b,\mathcal O_{X_b}).
\tag{9.18}
\]
On a smooth fibre its image is the complex line generated by the
\((0,1)\)-image of \(\gamma\).
\end{corollary}

\begin{proof}
The Leray spectral sequence and \(H^1(B,\mathcal O_B)=0\) identify
\[
 H^1(X,\mathcal O_X)=H^0(B,R^1f_*\mathcal O_X)=\mathbb C
\]
by \cref{thm:structure-direct-images-audit}.  Since the verified integral
topology gives \(H^1(X;\mathbb Z)=H^2(X;\mathbb Z)=0\), the analytic
exponential sequence
\[
0\to\mathbb Z\to\mathcal O_X
 \xrightarrow{u\mapsto e^{2\pi iu}}\mathcal O_X^*\to1
\]
identifies \(H^1(X,\mathcal O_X)\) with \(\operatorname{Pic}(X)\) and
shows that every line bundle is topologically trivial.  Under Leray,
restriction is evaluation of a global section of
\(R^1f_*\mathcal O_X\), followed by base change.  The only global section
is a scalar multiple of the nowhere-zero section \(\gamma^{0,1}\), so
evaluation is injective at every point and has the stated image on a
smooth fibre.
\end{proof}

\subsection{One line bundle non-torsion on every fibre}

\begin{theorem}[uniform non-torsion]
\label{thm:uniform-nontorsion-audit}
\Verified\ There is a line bundle \(L\in\operatorname{Pic}(X)\) such that
\(L|_{X_b}\) is non-torsion for every scheme-theoretic fibre
\(X_b=f^*(b)\), and \(L|_{S_j}\) is non-torsion for \(j=1,2\).
More precisely, identify \(\operatorname{Pic}(X)\) with \(\mathbb C\),
\(z\mapsto L_z\).  Every \(z\) outside the union of \(\mathbb R\) with a
countable set has these properties.
\end{theorem}

\begin{proof}
The pullback \(f^*\mathcal O_B(1)\) is nontrivial: by the projection
formula and \(f_*\mathcal O_X=\mathcal O_B\),
\[
 h^0(X,f^*\mathcal O_B(1))=h^0(B,\mathcal O_B(1))=2,
\]
whereas \(h^0(X,\mathcal O_X)=1\).  Rescale the additive coordinate in
\cref{cor:picard-line-audit} so that
\[
 L_{12}=f^*\mathcal O_B(1).
\tag{9.19}
\]

Fix a smooth fibre \(F_q\).  Put
\[
 \Lambda_q=\iota_q^{-1}\!\left(
 \operatorname{im}[H^1(F_q;\mathbb Z)\to
 H^1(F_q,\mathcal O_{F_q})]\right)\subset\mathbb C.
\tag{9.20}
\]
The integral image is a lattice in the underlying real vector space of
\(H^1(F_q,\mathcal O_{F_q})\), so \(\Lambda_q\) is discrete.  It contains
\(12\), because (9.19) is trivial on every fibre.  The exponential
sequence on \(F_q\) shows
\[
 L_z|_{F_q}\text{ is torsion}
 \quad\Longleftrightarrow\quad z\in\mathbb Q\Lambda_q.
\tag{9.21}
\]
If \(\operatorname{rk}\Lambda_q=1\), its nonzero element \(12\) forces
\(\Lambda_q\subset\mathbb R\); all its bad parameters are real.

Suppose \(\operatorname{rk}\Lambda_q=2\).  By
\cref{cor:picard-line-audit}, the image of \(\iota_q\) is the line
generated by the image of \(\gamma\).  Thus there is an integral class
\[
 \lambda=a\gamma+bu+cw+d\delta\notin\mathbb Q\gamma
\tag{9.22}
\]
whose \((0,1)\)-image lies on that line.  On the covering family, the four
marked lattice columns are
\[
 \pi_{\widehat\gamma}=(6\mu,\beta),\quad
 \pi_{\widehat u}=(\tau,\mu),\quad
 \pi_{\widehat w}=(1,0),\quad
 \pi_{\widehat\delta}=(0,1).
\tag{9.23}
\]
Every real cohomology class has a unique expression
\[
 2\operatorname{Re}(\bar\xi\,d\zeta_1+
                       \bar\upsilon\,d\zeta_2).
\]
Evaluating it on (9.23) gives
\[
\begin{aligned}
2\operatorname{Re}\bar\xi&=c,&
2\operatorname{Re}\bar\upsilon&=d,\\
2\operatorname{Re}(\bar\xi\tau+\bar\upsilon\mu)&=b,&
2\operatorname{Re}(6\bar\xi\mu+\bar\upsilon\beta)&=a.
\end{aligned}
\tag{9.24}
\]
For \(\gamma\), the first two right sides vanish.  The third equation then
shows that its \((0,1)\)-line is represented, up to a nonzero scalar, by
\((\operatorname{Im}\mu,-\operatorname{Im}\tau)\); the fourth equation
only fixes the scalar.  Write
\(\bar\xi=c/2+ir\), \(\bar\upsilon=d/2+is\).  The condition that the class
in (9.22) have \((0,1)\)-image on the \(\gamma\)-line is
\[
 \bar\xi\operatorname{Im}\tau+
 \bar\upsilon\operatorname{Im}\mu=0.
\]
Its real part, together with the third equation of (9.24), is equivalent
to the single holomorphic equation
\[
 c\tau+d\mu-b=0.
\tag{9.25}
\]
Neither \(a\) nor \(\beta\) occurs, as adding a multiple of \(\gamma\)
must not change the condition.

The functions \(1,\tau,\mu\) are linearly independent.  Applying the cusp
transformation \((\tau,\mu)\mapsto(\tau-1,\mu)\) to a constant linear
relation first kills the coefficient of \(\tau\).  If the remaining
relation made \(\mu\) constant, the order-four law
\(\mu\circ g_2=1+\mu/\tau\), with nonconstant \(\tau\), would force an
impossible identity.  Hence all coefficients vanish.

For every integral \(\lambda\notin\mathbb Q\gamma\), equation (9.25) is
therefore a nonzero holomorphic function on the covering Riemann surface.
Its zero set is discrete and countable.  There are countably many
\(\lambda\), so their union, and its image in \(B^\circ\), are countable.
This image contains every \(q\) for which
\(\operatorname{rk}\Lambda_q=2\).  Each such fibre contributes the
countable bad set \(\mathbb Q\Lambda_q\); all rank-one fibres together
contribute only parameters in \(\mathbb R\).  Thus the smooth fibres
exclude only \(\mathbb R\) together with a countable set.

There remain the three scheme-theoretic special fibres.  By
\cref{thm:structure-direct-images-audit,cor:picard-line-audit}, restriction
\[
 \mathbb C\longrightarrow H^1(X_b,\mathcal O_{X_b})
\]
is complex-linear and injective for each of them.  The exponential sequence
and countability of \(H^1(X_b;\mathbb Z)\) show that the parameters for
which \(L_z|_{X_b}\) is torsion form a countable set.

Finally consider the reductions \(S_j\).  The restriction
\(\varphi_j:\mathbb C\to H^1(S_j,\mathcal O_{S_j})\) is complex-linear.
It is nonzero.  Write \(\mathcal O_X(S_j)=L_{z_j}\).  Since
\[
 \mathcal O_X(m_jS_j)=f^*\mathcal O_B(1)=L_{12},
\]
the additive parameter gives \(z_j=12/m_j\), while
\cref{lem:finite-normal-order-audit} says that
\[
 L_{12/m_j}|_{S_j}=\mathcal O_X(S_j)|_{S_j}
\]
is nontrivial.  If \(\varphi_j\) were zero, the exponential sequence and
the vanishing of the restricted first Chern class would make every
\(L_z|_{S_j}\) trivial, a contradiction.  Thus \(\varphi_j\) is
injective, and the same exponential-sequence argument makes its torsion
parameters countable.  Removing these two additional countable sets proves
the theorem.
\end{proof}

\begin{remark}[provenance of the repair]
The period equation \(c\tau+d\mu-b=0\) and the countability strategy are
the candidate's.  The exact normal-character proof above makes the order of
\(\mathcal O_X(S_j)|_{S_j}\) explicit, while the cusp extension in
\cref{thm:structure-direct-images-audit} replaces the candidate's
collapse-dependent argument by the independently verified intrinsic
nearby-cycle map.  Identifying the restriction line with the global
\(\gamma^{0,1}\)-summand also removes the candidate's separate continuity
argument for a moving integral lift.
\end{remark}

\subsection{A local analytic obstruction to a cusp-only Picard argument}
\label{subsec:local-cusp-picard-obstruction}

The relative Picard space of the cusp family cannot by itself force the
normalised central restriction to vanish.  The obstruction is visible
directly in the toric model of \cref{thm:candidate-cusp-geometry}.  We give
the complete factor-of-automorphy calculation, including the descent
cocycle and the sign of the generic Chern class.

Retain the cusp notation of
\cref{prop:period-input-audit,thm:candidate-cusp-geometry}:
\[
 B_0=\begin{pmatrix}0&1\\-1&0\end{pmatrix},
 \qquad
 Z(s)=sB_0+C(t),
 \qquad
 C(t)=\begin{pmatrix}6\mu(t)&h(t)\\ b(t)-h(t)&\mu(t)\end{pmatrix},
 \qquad t=e^{2\pi i s},
\tag{5.17t}
\]
\[
 c_{\lambda}(t)=\exp(2\pi iC(t)\lambda),
 \qquad
 \exp(2\pi iZ(s)\lambda)=t^{B_0\lambda}c_\lambda(t)
 \quad(\lambda\in\bar\Lambda\simeq\mathbb Z^2),
\tag{5.17u}
\]
and
\[
 \Phi_\lambda(x,t)=(t^{B_0\lambda}x,t),
 \qquad
 \Psi_\lambda(p)=(c_\lambda(t(p)),1)\cdot\Phi_\lambda(p),
 \qquad
 \Psi_\lambda\Psi_\mu=\Psi_{\lambda+\mu}.
\tag{5.17v}
\]
Write \(E_v\), \(v\in\mathbb Z^2\), for the toric prime divisors of
\(Y_\epsilon\), so that \(E_0\simeq\mathrm{dP}_6\), and let
\[
 \pi:Y_\epsilon\longrightarrow N_0=Y_\epsilon/\bar\Lambda,
 \qquad
 \nu=\pi|_{E_0}:E_0\longrightarrow W
\]
be the quotient and the normalization of the central fibre.
Put
\[
 N'=\mathbb Z^2\oplus\mathbb Z,
 \qquad
 M'=\operatorname{Hom}(N',\mathbb Z).
\]
For \(m=(\ell,k)\in M'=\mathbb Z^2\oplus\mathbb Z\), fix the character
convention
\[
 \chi^m(x,t)=x^\ell t^k,
 \qquad
 \operatorname{div}(\chi^m)=
 \sum_{v\in\mathbb Z^2}\langle m,(v,1)\rangle E_v.
\tag{5.17w}
\]

\begin{proposition}[an explicit nonzero local cusp restriction]
\label{prop:local-cusp-picard-obstruction}
\Verified\ Let
\[
 H=\begin{pmatrix}1&0\\0&-6\end{pmatrix},
 \qquad
 P=B_0^{\mathsf t}H=\begin{pmatrix}0&6\\1&0\end{pmatrix},
 \qquad
 Q(v)=v^{\mathsf t}Hv=v_1^2-6v_2^2.
\tag{5.17x}
\]
Then
\[
 \mathcal D_Q=\sum_{v\in\mathbb Z^2}Q(v)E_v
\tag{5.17y}
\]
is a locally finite Cartier divisor on \(Y_\epsilon\).
The line bundle \(\mathcal O_{Y_\epsilon}(\mathcal D_Q)\) has an explicit
\((\bar\Lambda,\Psi)\)-equivariant structure and descends to a holomorphic
line bundle \(\mathcal L_Q\) on \(N_0\).  If
\[
 (a_1,a_2,a_3,a_4,a_5,a_6)
 =(e_1,e_2,e_2-e_1,-e_1,-e_2,e_1-e_2),
 \qquad C_i=E_0\cap E_{a_i},
\]
then
\[
 \nu^*(\mathcal L_Q|_W)
 \simeq
 \mathcal O_{E_0}\!\left(\sum_{i=1}^6Q(a_i)C_i\right),
\qquad
 \left(\deg_{C_i}\nu^*(\mathcal L_Q|_W)\right)_{i=1}^6
 =(-12,2,0,-12,2,0).
\tag{5.17z}
\]
In particular \(\nu^*(\mathcal L_Q|_W)\) is nontrivial.

On every nonzero fibre
\[
 F_t=\mathbb C^2/(Z(s)\mathbb Z^2+\mathbb Z^2),
\]
the first Chern class of \(\mathcal L_Q|_{F_t}\), in the ordered lattice
basis
\((\widehat\gamma,\widehat u,\widehat w,\widehat\delta)\) indexing the
columns of \([Z\mid I_2]\), is exactly
\[
 c_1(\mathcal L_Q|_{F_t})=-2q,
 \qquad
 q=u\wedge w+6\gamma\wedge\delta.
\tag{5.17aa}
\]
\end{proposition}

\begin{proof}
\emph{Local finiteness and Cartier status.}
The fan of \(Y\) consists of the cones over the cells of the locally finite
\(A_2\)-triangulation at height one.  If
\(U_\sigma\simeq\mathbb C^3\) is the affine chart belonging to a triangle
\(\sigma=\{v_0,v_1,v_2\}\), then \(U_\sigma\) meets precisely
\(E_{v_0},E_{v_1},E_{v_2}\).  Indeed, by the orbit--cone
correspondence, \(U_\sigma\cap E_v\ne\varnothing\) exactly when
\(\mathbb R_{\geq0}(v,1)\) is a face of the cone over \(\sigma\).
Thus \(\{E_v\}_{v\in\mathbb Z^2}\) is locally finite and
\(\mathcal D_Q\) is a well-defined locally finite Weil divisor.

Every triangle of the fan is unimodular, as proved in
\cref{thm:candidate-cusp-geometry}.  If \(z_0,z_1,z_2\) are the
characters dual to
\((v_0,1),(v_1,1),(v_2,1)\), then
\[
 \mathcal D_Q|_{U_\sigma}
 =\operatorname{div}\!
   \left(z_0^{Q(v_0)}z_1^{Q(v_1)}z_2^{Q(v_2)}\right).
\tag{5.17ab}
\]
Negative exponents merely make the displayed monomial meromorphic, as
permitted for a local equation of a Cartier divisor.  Hence
\(\mathcal D_Q\) is Cartier on every chart and therefore on
\(Y_\epsilon\).

\emph{The signed divisor pullback identity.}
For \(\lambda\in\bar\Lambda=\mathbb Z^2\), set
\[
 b_\lambda=B_0\lambda,
 \qquad
 \ell_\lambda=2Hb_\lambda,
 \qquad
 k_\lambda=Q(b_\lambda),
 \qquad
 m_\lambda=(\ell_\lambda,k_\lambda)\in M'.
\tag{5.17ac}
\]
The fan automorphism sends \(E_v\) to \(E_{v+b_\lambda}\).
Consequently \(\Phi_\lambda^*E_w=E_{w-b_\lambda}\), and
\begin{align*}
 \Phi_\lambda^*\mathcal D_Q-\mathcal D_Q
 &=\sum_{v\in\mathbb Z^2}
   \bigl(Q(v+b_\lambda)-Q(v)\bigr)E_v\\
 &=\sum_{v\in\mathbb Z^2}
   \bigl(2v^{\mathsf t}Hb_\lambda+
         b_\lambda^{\mathsf t}Hb_\lambda\bigr)E_v
 =\operatorname{div}(\chi^{m_\lambda}).
\end{align*}
The \(t\)-dependent horizontal torus translation
\(p\mapsto(c_\lambda(t(p)),1)\cdot p\) preserves every
torus-invariant prime divisor \(E_v\).  Hence
\[
 \boxed{\ 
 \Psi_\lambda^*\mathcal D_Q
 =\mathcal D_Q+\operatorname{div}(\chi^{m_\lambda})\ }.
\tag{5.17ad}
\]

\emph{The complete \(\Psi\)-cocycle.}
The matrices in (5.17t) and (5.17x) give
\[
 PC(t)=
 \begin{pmatrix}
 6(b(t)-h(t))&6\mu(t)\\
 6\mu(t)&h(t)
 \end{pmatrix},
\tag{5.17ae}
\]
which is symmetric.  Thus
\[
 A(t):=2PC(t)
\tag{5.17af}
\]
is a symmetric holomorphic matrix on the full cusp disc.  Define the
holomorphic units \(h_\lambda\) and meromorphic functions \(f_\lambda\) by
\[
 h_\lambda(t)=
 \exp\!\left(\pi i\lambda^{\mathsf t}A(t)\lambda\right),
 \qquad
 f_\lambda=h_\lambda(t)\chi^{m_\lambda}.
\tag{5.17ag}
\]
Equation (5.17ad) gives
\[
 \operatorname{div}(f_\lambda)
 =\Psi_\lambda^*\mathcal D_Q-\mathcal D_Q.
\tag{5.17ah}
\]

The cocharacter map underlying \(\Phi_\mu\) is
\[
 \varphi_\mu(y,y_3)=(y+y_3b_\mu,y_3),
\]
so its pullback on characters is
\[
 \varphi_\mu^*(\ell,k)
 =(\ell,k+\ell^{\mathsf t}b_\mu).
\tag{5.17ai}
\]
The quadratic identity for \(Q\) gives
\[
 m_{\lambda+\mu}=m_\mu+\varphi_\mu^*m_\lambda.
\tag{5.17aj}
\]
Indeed, the first coordinates add, while the last coordinate on the
right is
\[
 Q(b_\mu)+Q(b_\lambda)+2b_\lambda^{\mathsf t}Hb_\mu
 =Q(b_\lambda+b_\mu).
\]
Using (5.17u), (5.17v), and (5.17ai), one obtains
\begin{align*}
 \Psi_\mu^*\chi^{m_\lambda}
 &=\chi^{m_\lambda}(c_\mu(t),1)
   \chi^{\varphi_\mu^*m_\lambda}\\
 &=\exp\!\left(2\pi i\ell_\lambda^{\mathsf t}C(t)\mu\right)
   \chi^{\varphi_\mu^*m_\lambda}\\
 &=\exp\!\left(2\pi i\lambda^{\mathsf t}A(t)\mu\right)
   \chi^{\varphi_\mu^*m_\lambda}.
\tag{5.17ak}
\end{align*}
The last equality is exact because
\[
 \ell_\lambda^{\mathsf t}C\mu
 =(2HB_0\lambda)^{\mathsf t}C\mu
 =2\lambda^{\mathsf t}B_0^{\mathsf t}HC\mu
 =2\lambda^{\mathsf t}PC\mu
 =\lambda^{\mathsf t}A\mu.
\]
Since \(A=A^{\mathsf t}\), (5.17ag), (5.17aj), and (5.17ak) imply
\begin{align*}
 f_\mu\,\Psi_\mu^*f_\lambda
 &=h_\mu h_\lambda
   \exp(2\pi i\lambda^{\mathsf t}A\mu)
   \chi^{m_\mu+\varphi_\mu^*m_\lambda}\\
 &=\exp\!\left(
   \pi i(\lambda+\mu)^{\mathsf t}A(\lambda+\mu)\right)
   \chi^{m_{\lambda+\mu}}
 =f_{\lambda+\mu}.
\tag{5.17al}
\end{align*}

Use the divisor-sheaf convention
\[
 \mathcal O(D)(U)=
 \{g\in\mathcal M_Y(U):
   \operatorname{div}(g)+D|_U\geq0\}.
\]
If \(D'=D+\operatorname{div}(f)\), multiplication by \(f\) maps
\(\mathcal O(D')\) isomorphically to \(\mathcal O(D)\).  Thus
(5.17ad) and (5.17ag) define
\[
 \theta_\lambda:
 \Psi_\lambda^*\mathcal O(\mathcal D_Q)
 \xrightarrow{\ \times f_\lambda\ }
 \mathcal O(\mathcal D_Q).
\tag{5.17am}
\]
Because
\((\Psi_\lambda\Psi_\mu)^*=\Psi_\mu^*\Psi_\lambda^*\),
equation (5.17al) is exactly
\[
 \theta_{\lambda+\mu}
 =\theta_\mu\circ\Psi_\mu^*\theta_\lambda.
\tag{5.17an}
\]
This proves a complete equivariance datum, not merely invariance of the
divisor class.

\emph{Analytic descent.}
By \cref{thm:candidate-cusp-geometry}, the action on \(Y_\epsilon\) is
free and properly discontinuous.  Every point therefore has an open
neighbourhood \(U\) whose distinct translates are disjoint and on which
\(\pi|_U\) is a biholomorphism onto \(\pi(U)\).  Transport
\(\mathcal O(\mathcal D_Q)|_U\) to \(\pi(U)\), and on overlaps arising
from translates glue by (5.17am).  Equation (5.17an) is precisely the
triple-overlap condition.  The resulting holomorphic line bundle
\(\mathcal L_Q\) on \(N_0\) satisfies
\[
 \pi^*\mathcal L_Q
 \simeq\mathcal O_{Y_\epsilon}(\mathcal D_Q).
\tag{5.17ao}
\]

\emph{The normalized central restriction.}
Let \(j_0:E_0\hookrightarrow Y_\epsilon\).  Since
\(\nu=\pi\circ j_0\), (5.17ao) gives
\[
 \nu^*(\mathcal L_Q|_W)
 =j_0^*\pi^*\mathcal L_Q
 \simeq\mathcal O_{E_0}(\mathcal D_Q|_{E_0}).
\tag{5.17ap}
\]
The only \(E_v\), \(v\ne0\), meeting \(E_0\) are those indexed by the
six \(a_i\) in the statement, and the intersections are the six toric
\((-1)\)-curves \(C_i\).  Since \(Q(0)=0\), there is no
\(E_0|_{E_0}\) term, so
\[
 \mathcal D_Q|_{E_0}
 =\sum_{i=1}^6Q(a_i)C_i,
 \qquad
 (Q(a_1),\ldots,Q(a_6))=(1,-6,-5,1,-6,-5).
\tag{5.17aq}
\]
The hexagonal intersection table is
\[
 C_i^2=-1,
 \qquad
 C_i\cdot C_{i\pm1}=1,
 \qquad
 C_i\cdot C_j=0
 \quad(j\notin\{i,i\pm1\}),
\tag{5.17ar}
\]
with indices modulo six.  Therefore
\[
 \deg_{C_i}\mathcal O_{E_0}
 \!\left(\sum_jQ(a_j)C_j\right)
 =-Q(a_i)+Q(a_{i-1})+Q(a_{i+1}),
\]
and direct substitution gives
\[
 \begin{array}{c|rrrrrr}
 i&1&2&3&4&5&6\\ \hline
 -Q(a_i)+Q(a_{i-1})+Q(a_{i+1})
 &-12&2&0&-12&2&0.
 \end{array}
\tag{5.17as}
\]
The degree on \(C_1=C_{e_1}\) is \(-12\), proving nontriviality.

\emph{The generic-fibre factor and its exact Chern class.}
Fix \(t\ne0\), choose \(s\) with \(t=e^{2\pi i s}\), and write
\(x=e^{2\pi i\zeta}\), \(\zeta\in\mathbb C^2\).
Equation (5.17u) identifies the deck transformation indexed by
\(\lambda\in\mathbb Z^2\) with
\[
 \zeta\longmapsto\zeta+Z(s)\lambda
\]
modulo the horizontal translations
\(\zeta\mapsto\zeta+n\), \(n\in\mathbb Z^2\).  The matrix identities
\[
 HB_0=P^{\mathsf t},
 \qquad
 B_0^{\mathsf t}HB_0=PB_0
\tag{5.17at}
\]
give
\[
 \chi^{m_\lambda}(x,t)
 =x^{2P^{\mathsf t}\lambda}
  t^{\lambda^{\mathsf t}PB_0\lambda},
\qquad
 h_\lambda(t)
 =\exp(2\pi i\lambda^{\mathsf t}PC(t)\lambda).
\]
Since \(Z=sB_0+C\), it follows that
\[
 f_\lambda(\zeta)
 =\exp\!\left(
 2\pi i\lambda^{\mathsf t}PZ\lambda+
 4\pi i\lambda^{\mathsf t}P\zeta\right).
\tag{5.17au}
\]
The matrix \(PZ=sPB_0+PC\) is symmetric:
\[
 PB_0=\operatorname{diag}(-6,1),
\]
and \(PC\) is symmetric by (5.17ae).

Under the descent convention (5.17am), a pulled-back local section
written as \(\varphi(\zeta)\) in the canonical frame satisfies
\[
 \varphi(\zeta+Z\lambda)
 =f_\lambda(\zeta)^{-1}\varphi(\zeta).
\]
Thus the factor of automorphy for the lattice translation
\(\zeta\mapsto\zeta+Z\lambda+n\) is
\[
 a_{(\lambda,n)}(\zeta)
 =\exp(2\pi i g_{(\lambda,n)}(\zeta)),
\qquad
 g_{(\lambda,n)}(\zeta)
 =-\lambda^{\mathsf t}PZ\lambda
  -2\lambda^{\mathsf t}P\zeta.
\tag{5.17av}
\]
This is independent of \(n\) as a factor of automorphy: replacing
\(\zeta\) by \(\zeta+n\) changes \(g\) by the integer
\(-2\lambda^{\mathsf t}Pn\).

For \(l=(\lambda,n)\), \(l'=(\mu,m)\), the exponential connecting
cocycle is
\begin{align*}
 c(l,l')
 &=g_l(\zeta+Z\mu+m)+g_{l'}(\zeta)-g_{l+l'}(\zeta)\\
 &=-2\lambda^{\mathsf t}Pm\in\mathbb Z,
\tag{5.17aw}
\end{align*}
where the symmetry of \(PZ\) cancels the two mixed \(Z\)-terms.
The corresponding alternating integral form is
\[
 E(l,l')=c(l,l')-c(l',l)
 =-2\lambda^{\mathsf t}Pm+2\mu^{\mathsf t}Pn.
\tag{5.17ax}
\]
This is \(c_1(\mathcal L_Q|_{F_t})\).  In the ordered column basis of
\([Z\mid I_2]\), the cross-block of the class
\(q=u\wedge w+6\gamma\wedge\delta\) from
\((\widehat\gamma,\widehat u)\) to
\((\widehat w,\widehat\delta)\) is
\[
 \begin{pmatrix}
 q(\widehat\gamma,\widehat w)&
 q(\widehat\gamma,\widehat\delta)\\
 q(\widehat u,\widehat w)&
 q(\widehat u,\widehat\delta)
 \end{pmatrix}
 =\begin{pmatrix}0&6\\1&0\end{pmatrix}=P.
\tag{5.17ay}
\]
Hence (5.17ax) is exactly \(-2q\), as asserted.
\end{proof}

\begin{remark}[the exact obstruction to the purely analytic route]
\label{rem:local-cusp-picard-obstruction}
The proposition proves that the local homomorphism
\[
 \operatorname{Pic}(N_0)\longrightarrow\operatorname{Pic}(E_0),
 \qquad
 \mathcal L\longmapsto\nu^*(\mathcal L|_W),
\]
is not zero.  Thus the toric cusp geometry, conductor descent on \(W\),
or local relative-Picard representability cannot alone prove the desired
global vanishing
\[
 \nu^*i^*:\operatorname{Pic}(X)\longrightarrow
 \operatorname{Pic}(E_0).
\]
There is no contradiction with
\cref{prop:leray-edge-restriction,prop:direct-boundary-picard}: the local
bundle above has nonzero generic class \(-2q\), whereas the two global
proofs use the order-three and order-four fillings and their common
boundary relation to exclude every nonzero multiple of \(q\).

More generally, Bingener proves that if a holomorphic map is proper,
flat, and cohomologically flat in dimension zero, then its relative
Picard functor is represented by a \emph{not necessarily separated}
complex space
\cite[Introduction, p.~318; Satz~6.2, pp.~336--337]{Bingener1980}.
The same introduction explains that the exponential-sequence quotient
works directly under topological local triviality and cohomological
flatness in dimension one, but that passage to the quotient fails in
general.  The cusp map of \cref{thm:candidate-cusp-geometry} is not
topologically locally trivial across \(0\).  Bingener's theorem therefore
supplies representability but neither separatedness nor the
component-specialization assertion needed here.

Consequently a purely analytic completion must add global compatibility
with both finite fillings and must control the boundary support function;
neither relative-Picard representability nor conductor descent supplies
those statements.  The later calculation in
\cref{prop:analytic-ah-boundary,lem:analytic-support-specialization,%
thm:analytic-all-picard} now supplies them: it first forces the generic
Appell--Humbert form to vanish using the two finite fillings, and then
proves that its mixed block is the polarization of the three normalized
central degrees.  Thus this proposition remains the sharp cusp-local
obstruction, while the global purely analytic route is complete.
\end{remark}

\subsection{The canonical divisor and the exact anticanonical ring}
\label{sec:canonical-ring-recovery}

The analytic construction admits a further explicit calculation which does
not use its claimed identification with a sphere. Retain the actual proper
map $f:X\to B=\mathbb P^1$, the cusp $p_0$, the finite points $p_1,p_2$,
the reduced finite fibres $S_1,S_2$, and
$f^*(p_1)=3S_1$, $f^*(p_2)=4S_2$. All line-bundle isomorphisms below
refer to this constructed analytic space. No replacement period matrix or
new compactification is introduced.

\begin{theorem}[Canonical bundle and recovery of the original fibration]
\label{thm:canonical-ring-recovery}
The evaluation map
\[
 f^*(f_*\omega_{X/B})\longrightarrow\omega_{X/B}
\]
has zero divisor exactly $2S_2$. Consequently
\[
 \omega_X\simeq f^*\mathcal O_B(-p_2)\otimes\mathcal O_X(2S_2)
       \simeq\mathcal O_X(-2S_2),\qquad
 \omega_X^{-2}\simeq f^*\mathcal O_B(p_2).
 \tag{CR1}
\]
For every integer $k$ the actual direct image of the vertical divisor is
\[
 f_*\mathcal O_X(kS_2)
       =\mathcal O_B\bigl(\lfloor k/4\rfloor p_2\bigr)
 \tag{CR2}
\]
as subsheaves of the meromorphic-function sheaf, using pullback to identify
their sections. In particular, for each integer $m\geq0$,
\[
 h^0(X,\omega_X^{-m})=\lfloor m/2\rfloor+1,
 \qquad h^0(X,\omega_X^m)=0\quad(m>0).
 \tag{CR3}
\]
There is a graded algebra isomorphism
\[
 \mathbb C[U,V]\xrightarrow{\sim}
       \bigoplus_{m\geq0}H^0(X,\omega_X^{-m}),
       \qquad\deg U=1,\quad\deg V=2.
 \tag{CR4}
\]
The degree-two complete linear system has no base points and its map is
the original $f$, followed by a specified projective-coordinate map of $B$.
The degree-one system consists of the single divisor $2S_2$.
\end{theorem}

\begin{proof}
Write $\mathfrak K=f_*\omega_{X/B}$. The map is the natural adjunction
evaluation, not an identification inferred from the ranks of its two
terms. On a smooth torus fibre the invariant top form is nowhere zero.
At the cusp, the original form
$(dx_1/x_1)\wedge(dx_2/x_2)\wedge dt_c$ is a nowhere-zero canonical
frame in each of the unimodular toric charts. Its relative version is
the local generator of $\mathfrak K$ used in
\cref{lem:hodge-line-audit}, so evaluation has no zero there.

For a finite filling retain $t_j=s_j^{m_j}$,
$(m_1,m_2)=(3,4)$ and $(a_1,a_2)=(2,1)$. On its finite \'etale
total-space cover $q_j:J'_j\to N_j$, the actual direct-image generator
pulls back, by (9.5c)--(9.5f), to
\[
 s_j^{m_j-1-a_j}e^{\varphi_j(s_j)}
       (ds_j\wedge e)\otimes(dt_j)^{-1}.
 \tag{CR5}
\]
Here $(dt_j)^{-1}$ is the pulled-back inverse canonical frame on the
base, not division by $m_js_j^{m_j-1}ds_j$. The coefficient
$e^{\varphi_j}$ is a unit. Thus the zero divisor on the cover is
$(m_j-1-a_j)\{s_j=0\}$. Since $q_j$ is \'etale, faithful pullback
of Cartier divisors makes the downstairs divisor
$(m_j-1-a_j)S_j$. The two coefficients are respectively $0$ and $2$.
These neighbourhoods and the smooth locus cover $X$, proving the stated
global zero divisor with no omitted component.

The evaluation therefore identifies
\[
 \omega_{X/B}\simeq f^*\mathfrak K\otimes\mathcal O_X(2S_2).
\]
By \cref{lem:hodge-line-audit}, $\mathfrak K\simeq\mathcal O_B(p_0)$.
In the original affine base coordinate $t$, with
$t(p_1)=0$, $t(p_2)=1$, $t(p_0)=\infty$, retain the explicit choices
\[
 h=\frac1{t-1},\qquad \eta_B=\frac{dt}{(t-1)^2}.
\]
Their divisors are $p_0-p_2$ and $-2p_2$, respectively: at infinity,
with $w=1/t$, they are $w/(1-w)$ and $-dw/(1-w)^2$.
Thus $\eta_B$ specifies $\omega_B\simeq\mathcal O_B(-2p_2)$.
Tensoring with $f^*\omega_B$ gives
\[
 \omega_X\simeq
 \mathcal O_X\bigl(f^*(p_0)-2f^*(p_2)+2S_2\bigr).
\]
Multiplication by $f^*h$, with its displayed divisor, identifies the
base factor $\mathcal O_B(p_0-2p_2)$ with $\mathcal O_B(-p_2)$.
Using the actual multiplicity $f^*(p_2)=4S_2$ proves (CR1).

To prove (CR2), a section of $\mathcal O_X(kS_2)$ is represented by
a meromorphic function on $X$ with the indicated vertical order bound.
Away from $p_2$ it is holomorphic. Its restriction to each connected
compact fibre there is constant; $f_*\mathcal O_X=\mathcal O_B$
therefore descends it uniquely to a holomorphic function on that base
open set. Near $p_2$, pull back to the original cover $J'_2$, including
its zero section. The descended function becomes a function $a(s_2)$,
and $s_2^k a(s_2)$ is holomorphic at zero (with the evident vanishing
interpretation when $k<0$). Thus $a$ is meromorphic and has a Laurent
series bounded below. Invariance under the actual quotient, whose base
coordinate is $s_2\mapsto\zeta_2s_2$, says
$a(\zeta_2s_2)=a(s_2)$. Only exponents divisible by four occur, so
$a(s_2)=b(s_2^4)$ for a meromorphic $b$ on the original base disc.
The order inequality is exactly
\[
 4\operatorname{ord}_{p_2}(b)+k\geq0
 \quad\Longleftrightarrow\quad
 \operatorname{ord}_{p_2}(b)\geq-\lfloor k/4\rfloor.
\]
Conversely every such $b$ pulls back to an allowed meromorphic function.
This proves equality in both directions, as well as compatibility on
overlaps. No assumption about the field of all meromorphic functions
on $X$ is used. Applying (CR2) with $k=2m$ and $k=-2m$ and using
$h^0(\mathbb P^1,\mathcal O(d))=\max\{d+1,0\}$ proves (CR3).

Here is the exact algebra map in (CR4). Fix the isomorphism
$\omega_X^{-1}\simeq\mathcal O_X(2S_2)$ obtained above. Let $u$ be
the section corresponding to the canonical divisor section of
$\mathcal O_X(2S_2)$, so $\operatorname{div}(u)=2S_2$.
Its square, under the induced degree-two identification in (CR1),
is $f^*\sigma_{p_2}$, where $\sigma_{p_2}$ is the canonical divisor
section of $\mathcal O_B(p_2)$. Choose a second section $\ell$ of this
line bundle with $\ell(p_2)\ne0$, retaining that choice, and let $v$
correspond to $f^*\ell$ under the same degree-two identification.
Define $U\mapsto u$, $V\mapsto v$.

For $m=2r$, the monomials $u^{2r-2j}v^j$, $0\leq j\leq r$,
correspond exactly to
$f^*(\sigma_{p_2}^{r-j}\ell^j)$, a basis of
$H^0(B,\mathcal O_B(rp_2))$. For $m=2r+1$, multiplication by $u$
identifies the same basis with the full space: in the meromorphic
description (CR2), both $4rS_2$ and $(4r+2)S_2$ allow precisely
the base pole orders at most $r$. Thus the displayed map is an
isomorphism in every degree, proving both surjectivity and absence of
polynomial relations. Replacing $\ell$ by $a\ell+b\sigma_{p_2}$,
$a\ne0$, changes the second generator by exactly
$v\mapsto av+bu^2$; no scalar is silently discarded.
Replacing the chosen isomorphism $\omega_X^{-1}\to\mathcal O_X(2S_2)$
by $c$ times itself changes the generators by
$(u,v)\mapsto(c^{-1}u,c^{-2}v)$, with $c\in\mathbb C^*$ retained.

The two degree-two generators have no common zero because
$\sigma_{p_2},\ell$ do not. Their map is
$x\mapsto[\sigma_{p_2}(f(x)):\ell(f(x))]$, which is exactly
$f$ followed by that chosen coordinate map $B\to\mathbb P^1$.
Explicitly, writing $\ell=A/(t-1)+B$ with $A\ne0$, it is
$t\mapsto[t-1:A+B(t-1)]$, including the values $[0:A]$ at $p_2$
and $[1:B]$ at $p_0$.
Degree one has basis $u$ and its divisor is $2S_2$, as claimed.
\end{proof}

In the already specified additive Picard coordinate (9.19), with
$L_{12}=f^*\mathcal O_B(1)$, the identities retain the original constants:
\[
 \mathcal O_X(S_1)=L_4,\qquad
 \mathcal O_X(S_2)=L_3,\qquad
 \omega_X=L_{-6},\qquad \omega_X^{-2}=L_{12}.
\]
This last expression uses the previously proved Picard identification;
the divisor, direct-image and ring proofs above do not depend on it.

\subsection{The retained period constant is an effective deformation parameter}
\label{sec:retained-period-deformation}

Keep the functions, lattice marking, base coordinate and local fillings of
\cref{sec:candidate-geometry}. In particular, fix the actual cusp-regular
solution $\beta^{\rm part}$, not a replacement with a cusp pole, and write
\[
 \beta_c=\beta^{\rm part}+c,\qquad
 \mathcal P=\{c\in\mathbb C:\operatorname{Im}c<-M\},
 \quad M=\max_{B\setminus\{p_0\}}D_{\rm part}.
 \tag{PD1}
\]
The maximum and the exact admissible domain were proved in
\cref{prop:period-exact-admissible-constants}. The following calculation
uses the analytic construction, not its integral cohomology or sphere
recognition. Its purpose is to determine whether the freedom in $c$
actually changes the constructed complex manifold.

\begin{theorem}[Exact period deformation]
\label{thm:retained-period-deformation}
For every open disc $Q$ with $\overline Q\subset\mathcal P$, the retained
construction gives a proper holomorphic submersion
\[
 \pi:\mathcal X_Q\longrightarrow Q,
 \qquad F:\mathcal X_Q\longrightarrow B\times Q,
 \qquad \pi=\operatorname{pr}_Q\circ F,
 \tag{PD2}
\]
whose fibre over $c$ is the analytically constructed threefold $X_c$ and
whose restriction there is the original $f_c:X_c\to B$.
The three distinguished values remain the original
$t(p_1)=0$, $t(p_2)=1$, $t(p_0)=\infty$.

Fix a regular base point and one of its original lifts $z$, and fix
$c_*\in Q$. With all functions below evaluated at that $z$, set
\[
 T=\operatorname{Im}\tau>0,\quad m=\operatorname{Im}\mu,\quad
 D=\operatorname{Im}\beta_{c_*}-\frac{6m^2}{T}<0,\quad a=\frac mT,
 \qquad P=\begin{pmatrix}0&0\\-a&1\end{pmatrix}.
 \tag{PD3}
\]
For $\delta=c-c_*$, the exact marked real-linear map from the universal
cover of the fibre at $c_*$ to that at $c$ is
\[
 w=A_\delta\zeta+B_\delta\overline\zeta,\qquad
 A_\delta=I+\frac{\delta}{2iD}P,\qquad
 B_\delta=-\frac{\delta}{2iD}P.
 \tag{PD4}
\]
It maps the original period lattice onto the new period lattice, and
therefore descends to a real-analytic marked diffeomorphism of the two
tori. In the explicit convention that the transported $(0,1)$ tangent
space is spanned by
$\partial_{\overline\zeta_k}+\sum_j\nu^j_k\partial_{\zeta_j}$,
its exact coefficient matrix and its first derivative are
\[
 \nu_\delta=\frac{\delta}{2iD+\delta}P,\qquad
 \dot\nu_z=\frac{1}{2iD}\,
 \partial_{\zeta_2}\otimes
 (d\overline\zeta_2-a\,d\overline\zeta_1).
 \tag{PD5}
\]
The class of $\dot\nu_z$ in
$H^{0,1}(f_{c_*}^{-1}(t(z)),T^{1,0})$ is nonzero.
With the tangent-exact-sequence convention
$\operatorname{KS}(\partial_c)=[\bar\partial h]$ for a smooth
$(1,0)$ lift $h$ of $\partial_c$, the actual fibre class is
$\operatorname{KS}_z(\partial_c)=-[\dot\nu_z]$.
Moreover the first-order deformation of the entire $X_{c_*}$ obtained
from (PD2) is nontrivial: its Kodaira--Spencer class in
$H^1(X_{c_*},T_{X_{c_*}})$ is nonzero. In particular
$h^1(X_{c_*},T_{X_{c_*}})\geq1$.
\end{theorem}

\begin{proof}
\emph{Construction with the parameter retained.}
On an ordinary lifted base open set use the original period map
\[
 \Pi(z,c)=
 \begin{pmatrix}
 6\mu(z)&\tau(z)&1&0\\
 \beta^{\rm part}(z)+c&\mu(z)&0&1
 \end{pmatrix}:
 \Lambda\otimes\mathbb R\longrightarrow\mathbb C^2.
 \tag{PD6}
\]
Its real determinant is $\operatorname{Im}\tau\,D_c$, which never
vanishes on $Q$. Thus $(z,c,u)\mapsto[z,c,\Pi(z,c)u]$ identifies the
quotient, as a real-analytic family, with the original marked real torus
over this base open set times $Q$. The lattice action is holomorphic in
$(z,c,\zeta)$; it is free and properly discontinuous locally uniformly
on compact base-parameter sets, since the inverse real period matrices
are uniformly bounded there. The quotient is therefore a complex
manifold. Projection to $Q$ is a submersion because it is so before
the covering quotient. Its map to the lifted base times $Q$ is proper:
over any compact subset the image of that subset times a fixed closed
lattice parallelotope is the full inverse image.

The deck matrices are the original integral $M_g$ and the original
matrices $R_g(z)$, with
$\Pi(gz,c)=R_g(z)\Pi(z,c)M_g$. The $R_g$ do not depend on $c$.
Consequently the ordinary base quotient is jointly holomorphic.
Proper discontinuity there follows from proper discontinuity on the
base, exactly as in \cref{lem:middle-period-proper}, with compact
parameter sets carried along.

At a finite filling retain $m_1=3$, $m_2=4$,
$s_j(g_jz)=e^{-2\pi i/m_j}s_j(z)$, and the actual logarithmic lift
\[
 (z,c,\zeta)\longmapsto
 \left(g_jz,c,R_{g_j}(z)\zeta+
                  \Pi(g_jz,c)\frac{v_j}{m_j}\right).
 \tag{PD7}
\]
In the real marking this is $(z,c,u)\mapsto
(g_jz,c,A_ju+v_j/m_j)$, independent of $c$.
Every nonidentity power remains free: the invariant integral covector
$\gamma$ takes the translation of its $k$-th power to $k/3$ or $-k/4$,
which is not an integer for $0<k<m_j$.
The finite quotient is therefore a local biholomorphism and is a
submersion to $Q$. The same compact-parallelotope argument proves
properness over $D_j\times Q$. Its map is still exactly
$t_j=s_j^{m_j}$. On the punctured overlap the translating section is
\[
 \frac{\log s_j(z)}{2\pi i}\Pi(z,c)v_j.
 \tag{PD8}
\]
A change of logarithm adds an actual period. It is therefore a
branch-independent holomorphic torus section, jointly in $z,c$.
The calculation (2.9e) applies with $c$ retained and proves the
punctured-overlap conjugacy, not merely equality of the monodromy orders.

For the cusp retain the same infinite unimodular fan and the full matrix
\[
 C(t_c,c)=
 \begin{pmatrix}6\mu(t_c)&h(t_c)\\
 b^{\rm part}(t_c)-h(t_c)+c&\mu(t_c)
 \end{pmatrix},\qquad
 B_0=\begin{pmatrix}0&1\\-1&0\end{pmatrix}.
 \tag{PD9}
\]
Choose a closed cusp disc of radius $r_1$ on which these functions are
holomorphic. Put
\[
 K_Q=\max_{|t_c|\le r_1,\ c\in\overline Q}
          \|-2\pi\operatorname{Im}C(t_c,c)\|,
 \qquad 0<\epsilon<\min(r_1,1),\quad
          |\log\epsilon|>2K_Q.
 \tag{PD10}
\]
The action on $Y_\epsilon\times Q$ is the original one,
\[
 (x,t_c,c)\longmapsto
 \left(e^{2\pi iC(t_c,c)\lambda}
                  t_c^{B_0\lambda}x,t_c,c\right),
 \quad\lambda\in\mathbb Z^2.
 \tag{PD11}
\]
It is holomorphic across the boundary by the same toric automorphism
and nowhere-zero torus multiplier. The estimates in (2.8) now hold
uniformly on $Q$:
\[
 B_{t_c,c}=B_0+
       \frac{-2\pi\operatorname{Im}C(t_c,c)}{\log|t_c|},\quad
 \|B_{t_c,c}-B_0\|<\tfrac12,\quad
 \|B_{t_c,c}\lambda\|>\tfrac12\|\lambda\|\quad(\lambda\ne0).
 \tag{PD12}
\]
They exclude stabilizers on the dense torus and bound the finite set
of translations that can join each pair of bounded charts, uniformly
in $c$. At $t_c=0$ a nonzero translation moves the bounded cell, so
there is no stabilizer there either. A compact set is contained in finitely
many of the open charts $\mathcal U_\sigma(2,\epsilon)\times Q$ from
\cref{lem:toric-chart-control}.
The chart-pair estimate proves proper discontinuity on the whole
product, including sequences tending to its boundary.

For clarity properness is also uniform, not only pointwise in $c$.
For $\epsilon_1<\epsilon$, let $K_0$ be the finite union of closed
unit polydiscs associated to triangles meeting the closed radius
$\sqrt2$ ball, intersected with $|t_c|\le\epsilon_1$, exactly as in
\cref{thm:candidate-cusp-geometry}. For every $c$ and $t_c\ne0$,
choose $\lambda$ so that
$\|B_{t_c,c}^{-1}y+\lambda\|_\infty\le1/2$.
The translated logarithmic position has norm at most
$3/(2\sqrt2)<\sqrt2$, and hence lies in a chart defining $K_0$.
At $t_c=0$, the six closed bidiscs covering $E_0$ give the same
representatives after the exact translation by $B_0v$, with multiplier
$\exp(2\pi iC(0,c)B_0v)$ retained. Therefore, for compact $K\subset Q$,
the image of the compact $K_0\times K$ is the full inverse image of
$\{|t_c|\le\epsilon_1\}\times K$. Every compact set in the cusp
base times $Q$ is closed in a set of this form. This proves properness.
The quotient is a submersion to $Q$ because its covering charts are
open subsets of $Y_\epsilon\times Q$.

Exponentiating the original logarithmic torus coordinates gives the
jointly holomorphic cusp overlap, with period identity
$Z=sB_0+C(t_c,c)$. Glue it and (PD8) to the middle family.
The special base discs are disjoint. Each resulting inverse image of
a member of this base cover is its entire Hausdorff local piece.
Two points with different $(t,c)$ values are separated on the base;
two with the same value lie in one Hausdorff open piece and can be
separated there. This proves Hausdorffness. Local complex charts give
a complex fourfold and the submersion $\pi$.
Properness of $F$ follows by covering any compact set of $B\times Q$
by finitely many relatively compact sets inside these base pieces and
using the just proved local properness. Since $B$ is compact, projection
$B\times Q\to Q$ is proper, proving (PD2). The fibres and overlaps
are precisely those in the original analytic construction.

\emph{The exact variation on a regular fibre.}
Write $u=(u_1,u_2,u_3,u_4)^t$ in the original ordered lattice basis and
$\zeta=\Pi(z,c_*)u$. Taking imaginary parts gives the two exact equations
\[
 \operatorname{Im}\zeta_1=6m u_1+T u_2,\qquad
 \operatorname{Im}\zeta_2=(D+6m^2/T)u_1+m u_2.
 \tag{PD13}
\]
Subtracting $a$ times the first equation from the second yields
\[
 u_1=\frac{\operatorname{Im}\zeta_2-a\operatorname{Im}\zeta_1}{D}.
 \tag{PD14}
\]
In the changed periods the same $u$ maps to
$w=\Pi(z,c)u=\zeta+\delta e_2u_1$. Substitution of
$\operatorname{Im}\zeta=(\zeta-\overline\zeta)/(2i)$ proves (PD4).
This also proves its asserted lattice map and real invertibility:
it is exactly $\Pi(z,c)\Pi(z,c_*)^{-1}$ as a real-linear map.
Its real determinant is $(D+\operatorname{Im}\delta)/D>0$.
Neither the nonholomorphic part nor the original marking is dropped.

Here $P^2=P$, and
\[
 \det A_\delta=1+\frac\delta{2iD},\qquad
 A_\delta^{-1}=I-\frac{\delta}{2iD+\delta}P.
 \tag{PD15}
\]
The denominator can vanish only for $\delta=-2iD$. At that parameter
$D+\operatorname{Im}\delta=-D>0$, which is outside the admissible
domain. Thus this inverse is defined at every $c\in Q$.
Annihilating $w$ by the declared new $(0,1)$ fields gives
$A_\delta\nu_\delta+B_\delta=0$, which proves (PD5), including
its sign. All its coefficients are constant in the fibre coordinates.
Thus its first derivative is a globally defined $\bar\partial$-closed
tangent-valued $(0,1)$ form on the torus.

It cannot be $\bar\partial V$ for a globally defined smooth vector
field $V$. Indeed the invariant holomorphic frame
$\partial_{\zeta_1},\partial_{\zeta_2}$ writes $V$ with periodic
smooth coefficients. The integral of each of their derivatives
$\partial_{\overline\zeta_k}V^j$ against the translation-invariant
probability volume is zero: real constant vector fields preserve this
volume, and integration of their derivatives on the compact torus is
zero. But the coefficient of
$\partial_{\zeta_2}\otimes d\overline\zeta_2$ in $\dot\nu_z$ is
the nonzero constant $1/(2iD)$. Its integral is that same constant.
This contradiction proves nonzero Dolbeault class without a dimension
count or an appeal to numerical period variation.

The conversion to the tangent-exact-sequence convention retains a
minus sign. In the moving coordinates the real-marking lift has
$(1,0)$ part $h=\partial_c+u_1\partial_{w_2}$ at $c_*$.
Since (PD14) gives
\[
 \bar\partial u_1=-\frac{d\overline\zeta_2-a\,d\overline\zeta_1}{2iD},
\]
the connecting class $[\bar\partial h]$ is exactly $-[\dot\nu_z]$.
Both conventions therefore give a nonzero class; they are not identified
without this sign conversion.

\emph{The intrinsic pencil in the parameter family.}
Let $\mathcal S_2\subset\mathcal X_Q$ be the reduced finite divisor
over $p_2\times Q$, so $F^*(p_2\times Q)=4\mathcal S_2$.
The proof of \cref{thm:canonical-ring-recovery} works for relative
forms over $Q$, as can be checked on each original chart.
On the middle family the frame is $d\zeta_1\wedge d\zeta_2$ relative
to $B\times Q$. At the cusp the relative canonical frame over $Q$ is
$(dx_1/x_1)\wedge(dx_2/x_2)\wedge dt_c$.
Extra $dc$ terms in transition differentials vanish in relative forms.
The finite determinant characters $\chi_j$, their correcting units
$e^{\varphi_j(s_j)}$, and the modular function $\Theta$ used in
(9.4)--(9.6) are independent of $c$, since $R_g,\tau,\mu$ are.
Thus their coefficient identities hold jointly over $Q$.
The evaluation exponents remain $m_j-1-a_j=0,2$.
It follows by the actual same transition functions and divisors that
\[
 \omega_{\mathcal X_Q/Q}\simeq\mathcal O_{\mathcal X_Q}(-2\mathcal S_2),
 \qquad
 \omega_{\mathcal X_Q/Q}^{-2}\simeq
          F^*\operatorname{pr}_B^*\mathcal O_B(p_2).
 \tag{PD16}
\]
Here the base identifications retain $h=1/(t-1)$ and
$\eta_B=dt/(t-1)^2$ from (CR1); these functions also do not depend on $c$.
In particular (PD16) is a relative line-bundle isomorphism, not a
collection of separately chosen isomorphisms of its fibres.

We need its complete degree-two system after an infinitesimal base
change, and prove this explicitly. Put
$A=\mathbb C[\epsilon]/(\epsilon^2)$, base-change by
$c=c_*+\epsilon$, write $X_A$ for the resulting deformation and
$L_A=\omega_{X_A/A}^{-2}$. For $L_*=\omega_{X_{c_*}}^{-2}$,
flatness of the deformation and local freeness of $L_A$ give
\[
 0\longrightarrow L_*\xrightarrow{\ \epsilon\ }L_A
        \longrightarrow L_*\longrightarrow0.
 \tag{PD17}
\]
The two sections of $\mathcal O_B(p_2)$ used in (CR4), namely
$\sigma_{p_2}$ and $\ell=A_0/(t-1)+B_0'$ with $A_0\ne0$,
pull back by (PD16) to sections $s_0,s_1$ of $L_A$ whose reductions
are a basis of $H^0(X_{c_*},L_*)$. This basis assertion is exactly
the proved degree-two case of (CR2)--(CR4).
Any section of $L_A$ reduces to a complex linear combination of these
two reductions; subtract that lifted combination. Exactness of (PD17)
then writes the remainder as $\epsilon$ times a section of $L_*$,
again a combination of that basis. Conversely, if an $A$-linear
combination of $s_0,s_1$ is zero, reduce modulo $\epsilon$ and then
use the injective first arrow of (PD17) to see that all four complex
coefficients vanish. Therefore
\[
 H^0(X_A,L_A)=A s_0\oplus A s_1.
 \tag{PD18}
\]
The system is basepoint-free and its morphism is $f_A$ followed by
the fixed coordinate map $t\mapsto[t-1:A_0+B_0'(t-1)]$.
This proves the required base-change statement without presuming that
pointwise dimensions alone imply it.

\emph{Nontriviality for the complete threefold.}
Suppose its first-order deformation were trivial, with an isomorphism
$X_A\simeq X_{c_*}\times\operatorname{Spec}A$ reducing to the identity.
Intrinsic relative canonical bundles and their complete spaces of
sections are preserved by that isomorphism. By (PD18), it would
therefore identify $f_A$ with $a_A\circ f_{c_*}$ for a projective
linear change $a_A\in\operatorname{PGL}_2(A)$ reducing to the identity.

This change fixes all three distinguished points even to first order.
Here is a scheme-sensitive check. In the finite covering charts
$f_A$ is $t_j=s_j^{m_j}$, so its relative critical ideal is
$(s_j^{m_j-1})$, and its restriction to that critical scheme has
$t_j=0$. The map from the parameter ring $A$ to this critical
scheme is injective; it contains a constant-parameter section given
by $s_j=0$ and the zero section of the covering torus.
At the cusp, in a chart containing a triple point,
$t_c=z_0z_1z_2$ and the critical ideal is
$(z_1z_2,z_0z_2,z_0z_1)$. Again $t_c=0$ in its ring and $A$ injects,
using $z_0=z_1=z_2=0$. At a double chart the same assertion follows
from $t_c=z_0z_1$; a smooth chart has no critical point.
These critical ideals are intrinsic to the differential of the map
over $A$, and the charts are local biholomorphisms onto $X_A$.
Away from the three distinguished base values the map is a
submersion, so it has no critical point.
Consequently the scheme-theoretic critical-value sections are exactly
the three constant sections $p_1,p_2,p_0$. An automorphism reducing
to the identity cannot permute these three distinct closed points.
Its value at each is therefore the same point over $A$, not a
nonzero infinitesimal displacement: such a displacement would be a
nonzero element of $\epsilon\mathbb C$ in an injective copy of $A$
on the just displayed critical chart.

For completeness write $a_A(t)=t+\epsilon(b_0+b_1t+b_2t^2)$.
Fixing $0$ gives $b_0=0$, fixing $1$ gives $b_0+b_1+b_2=0$,
and fixing infinity, in $u=1/t$, gives $b_2=0$.
Hence $a_A$ is the identity. The hypothesized trivialization thus
preserves $f$, and restricting it to any regular base value would
trivialize the torus deformation (PD4). Its nonzero class (PD5)
excludes this.

Finally, the zero Kodaira--Spencer class is precisely the condition for
such a first-order trivialization. One may see the needed direction
directly: choose local product charts for the holomorphic submersion
$\pi$ and differentiate their transition functions at $c_*$. The
resulting tangent-valued overlap functions form a Cech $1$-cocycle.
If its class vanishes, after refinement it is the coboundary of local
holomorphic vector fields. Changing each chart by that vector field
times $\epsilon$ cancels the first-order overlap terms, producing
the identity-reducing trivialization. This is the usual
Kodaira--Spencer connecting class (for the standard deformation-space
interpretation see \cite[Section~67, p.~15]{McMullen1996Complex});
the explicitly recorded conversion
from the plus-graph convention in (PD5) changes its sign, not its
vanishing. The contradiction proves the final assertion.
\end{proof}

The theorem establishes a genuinely varying analytic family with its
original marked fibration and its intrinsic degree-two canonical pencil.
It does not assert that every two distinct parameters give nonisomorphic
unmarked manifolds, or that this family exhausts their deformation space.
The global integral topology and the CDP argument are not inputs to this
deformation calculation.


The global identifications within this retained family are computed in
\cref{pshift:thm:classification}: they are precisely $c'-c\in2\mathbb Z$,
including the finite-filling obstruction and every cusp overlap.


% New proof fragment. No canonical source has been edited by this lane.
% Uses the actual period construction, finite normalization certificate,
% toric fan, and analytically proved canonical-pencil recovery.
\subsection{Exact global identifications of the retained period parameter}
\label{pshift:sec:classification}

Keep the original base $B$, coordinate $t$, values
$t(p_1)=0,t(p_2)=1,t(p_0)=\infty$, and the original functions
$\tau,\mu,\beta^{\mathrm{part}}$. Write
\[
 \beta_c=\beta^{\mathrm{part}}+c,\qquad
 \mathcal P=\{c\in\mathbb C:\operatorname{Im}c<-M\},\qquad
 \Pi_c(z)=\begin{pmatrix}6\mu(z)&\tau(z)&1&0\\
                      \beta_c(z)&\mu(z)&0&1\end{pmatrix}.
 \tag{PS1}\label{pshift:eq:periods}
\]
All assertions below concern the analytic pieces and their actual gluing.
No integral cohomology, sphere recognition, or CDP assertion is an input.
In particular, the global parameter classification is not inferred from
an equality between two fibre lattices.

Put $e=(0,1)^{\mathsf t}\in\mathbb C^2$ and retain the ordered lattice
\[
 \Lambda=\mathbb Z\langle\widehat\gamma,\widehat u,
                        \widehat w,\widehat\delta\rangle,
 \qquad E=\widehat\delta\otimes\gamma.
 \tag{PS2}\label{pshift:eq:rank-one-map}
\]
Thus $E$ is the matrix with its only nonzero entry, equal to one, in
row four and column one; $E^2=0$ and $\Pi_cE=e\gamma$.
The finite data, with every sign retained, are
\[
\begin{gathered}
 A_1=\begin{pmatrix}1&0&0&0\\6&0&1&0\\-6&-1&-1&0\\-2&1&0&1\end{pmatrix},
 \qquad
 A_2=\begin{pmatrix}1&0&0&0\\0&0&-1&0\\-6&1&0&0\\3&0&1&1\end{pmatrix},\\
 v_1=(1,2,-4,0)^{\mathsf t},\quad
 v_2=(-1,-3,3,0)^{\mathsf t},\quad
 (m_1,m_2)=(3,4),\quad
 (\epsilon_1,\epsilon_2)=(1,-1),\quad \epsilon_j=\gamma(v_j).
\end{gathered}
\tag{PS3}\label{pshift:eq:finite-data}
\]
The covariance is $R_g\Pi_c=\Pi_c(gz)A_g$, with
\[
 R_{g_1}=\begin{pmatrix}-1/\tau&0\\(1-\mu)/\tau&1\end{pmatrix},\qquad
 R_{g_2}=\begin{pmatrix}1/\tau&0\\-\mu/\tau&1\end{pmatrix},\qquad
 R_{g_0}=I_2.
 \tag{PS4}\label{pshift:eq:covariance}
\]
In particular $R_ge=e$ for every $g$.

\begin{theorem}[The exact identification subgroup is $2\mathbb Z$]
\label{pshift:thm:classification}
For $c,c'\in\mathcal P$, the retained analytic threefolds are
biholomorphic if and only if
\[
                         c'-c\in2\mathbb Z.
 \tag{PS5}\label{pshift:eq:classification}
\]
Every such biholomorphism preserves the original map to $B$. For every
even integer $n$, the proof below constructs an isomorphism
$X_c\longrightarrow X_{c+n}$ whose complex-linear part on the original
universal fibre coordinates is $I_2$. Its map on the marked fibre lattice
is $I-nE$.

The equality of the middle lattices for all integer shifts loses a
nonzero order-two finite-filling invariant: the shift $n$ changes that
invariant by $n\bmod2$. The cusp contributes no additional obstruction
to the displayed even-shift construction.
\end{theorem}

\subsubsection{The middle family and its complete rational centralizer}

For an integer $n$, put $S_n=I+nE$. Since $S_n^{-1}=S_{-n}$,
\[
 \Pi_{c+n}=\Pi_c S_n,\qquad
 \Pi_{c+n}\Lambda=\Pi_c\Lambda.
 \tag{PS6}\label{pshift:eq:integer-shear}
\]
Both original matrices fix $\widehat\delta$ and preserve $\gamma$.
Consequently $A_jE=EA_j=E$, and the same commutation holds for every
monodromy word. The identity of $\mathbb C^2$ therefore identifies
the two full middle torus families with their actual $R_g$-action.
Its source-to-target lattice map is $S_{-n}$, since
$\Pi_{c+n}S_{-n}=\Pi_c$.

\begin{lemma}[All middle-family linear parts]
\label{pshift:lem:centralizer}
The simultaneous rational centralizer of $A_1,A_2$ is exactly
\[
 \{aI+bE:a,b\in\mathbb Q\}.
 \tag{PS7}\label{pshift:eq:centralizer}
\]
A biholomorphism between the middle families for $c,c'$ over the
identity of $B^\circ$ forces $c'-c\in\mathbb Z$. Its complex-linear
part on every fibre is either $I_2$ or $-I_2$.
\end{lemma}

\begin{proof}
Solving the two fixed-vector systems from (PS3) gives
\[
 (\Lambda\otimes\mathbb Q)^{A_1}
   =\mathbb Q(1,2,-4,0)^{\mathsf t}+\mathbb Q\widehat\delta,
 \quad
 (\Lambda\otimes\mathbb Q)^{A_2}
   =\mathbb Q(1,3,-3,0)^{\mathsf t}+\mathbb Q\widehat\delta.
\]
Their intersection is $\mathbb Q\widehat\delta$. The fixed covectors
are respectively
\[
 \mathbb Q\gamma+\mathbb Q(2u+w+3\delta),\qquad
 \mathbb Q\gamma+\mathbb Q(u+w+2\delta).
\]
Their intersection is $\mathbb Q\gamma$: comparing the $u$ and $w$
coefficients of a common multiple gives $2x=y$ and $x=y$, hence
$x=y=0$. Every commuting endomorphism therefore preserves
$\mathbb Q\widehat\delta$ and $\ker\gamma$.

On $\ker\gamma/\mathbb Q\widehat\delta$, in the classes of
$(\widehat u,\widehat w)$, the two matrices are
\[
 Q_1=\begin{pmatrix}0&1\\-1&-1\end{pmatrix},\qquad
 Q_2=\begin{pmatrix}0&-1\\1&0\end{pmatrix}.
\]
A matrix commuting with $Q_2$ has the form
$\left(\begin{smallmatrix}b&d\\-d&b\end{smallmatrix}\right)$.
Commutation with $Q_1$ gives $d=0$. Thus a commuting endomorphism $H$
has columns
\[
 H\widehat\gamma=(a,r,s,k)^{\mathsf t},\quad
 H\widehat u=b\widehat u+u_0\widehat\delta,\quad
 H\widehat w=b\widehat w+w_0\widehat\delta,\quad
 H\widehat\delta=d_0\widehat\delta.
\]
Use the actual columns $A_1\widehat u=-\widehat w+\widehat\delta$,
$A_1\widehat w=\widehat u-\widehat w$, and
$A_2\widehat u=\widehat w$. Commutation on these columns gives
\[
 d_0-w_0=b+u_0,\qquad u_0=2w_0,\qquad w_0=u_0.
\]
Hence $u_0=w_0=0$ and $d_0=b$. Commutation on
$\widehat\gamma$ with $A_1$ gives
\[
 r=2(a-b),\qquad s=-4(a-b).
\]
Commutation on that column with $A_2$ also gives $r=-s$.
Therefore $a=b$, $r=s=0$, and $H=aI+kE$. Conversely these matrices
commute, by the fixed-vector and fixed-covector identities preceding
(PS6). This proves (PS7).

For the geometric assertion, a holomorphic map between two compact
complex tori lifts to their universal covers. The derivative of the
lift is periodic, so its entries descend to holomorphic functions on
the compact source torus and are constant. The lift is consequently
affine, with a complex-linear part $L$. In a holomorphic family these
linear parts vary holomorphically. Their lattice maps are integral
and locally constant, hence give an isomorphism of the two marked
local systems. Because the map is over the identity of the base,
that integral map $H$ commutes with both $A_j$. Thus
$H=aI+bE$, with $a,b\in\mathbb Z$. Its determinant is $a^4$;
invertibility over $\mathbb Z$ gives $a=1$ or $a=-1$.

The complete period equation is
\[
                   L(z)\Pi_c(z)=\Pi_{c'}(z)(aI+bE).
 \tag{PS8}\label{pshift:eq:period-morphism}
\]
Its third and fourth columns give $L(z)=aI_2$. Its first column
then gives $a(c'-c)+b=0$. Therefore $c'-c=-b/a\in\mathbb Z$,
and the asserted two linear parts exhaust the possibilities.
\end{proof}

\subsubsection{The exact finite obstruction, including normalization}

The two cyclic norm matrices are, by multiplication of (PS3),
\[
 N_1=I+A_1+A_1^2=
 \begin{pmatrix}3&0&0&0\\6&0&0&0\\-12&0&0&0\\0&2&1&3\end{pmatrix},
\]
\[
 N_2=I+A_2+A_2^2+A_2^3=
 \begin{pmatrix}4&0&0&0\\12&0&0&0\\-12&0&0&0\\0&2&2&4\end{pmatrix}.
 \tag{PS9}\label{pshift:eq:norms}
\]
For example these formulas give
$N_1\widehat w=\widehat\delta$,
$N_2\widehat u=N_2\widehat w=2\widehat\delta$,
$N_1\widehat\delta=3\widehat\delta$, and
$N_2\widehat\delta=4\widehat\delta$.

We first justify why a local isomorphism must satisfy the exact affine
equations on the original finite cover. The actual ramified base change
$t_j=u^{m_j}$ has local equation
\[
 s_j^{m_j}-u^{m_j}
   =\prod_{a=0}^{m_j-1}(s_j-e^{2\pi ia/m_j}u)=0.
\]
Its normalization separates all these smooth branches, each with
coordinates $(z_1,z_2,u)$. The original map
$y\mapsto(q_j(y),s_j(y))$ identifies the original torus family $J'_j$
with this normalization, as proved by the explicit branches in
\cref{ffcert:thm:normalization}. A biholomorphism over the identity of
$D_j$ base-changes over the identity of the $u$-disc and lifts uniquely
to the normalization: on local rings it carries integral elements in
the total fraction ring to integral elements, and its inverse does
the same. This lift commutes with the original deck action because
both composites lift the same map $(x,u)\mapsto(F(x),\zeta_j u)$.
In particular one may not replace the generator by an arbitrary
coprime power.

Let $c'=c+n$, where $n$ is an integer, and suppose the linear part of
the lifted map is $aI_2$, $a\in\{1,-1\}$. The lifted map on the torus
family is $[z,\zeta]\mapsto[z,a\zeta+B_j(z)]$. Its translating
section has a holomorphic lift $B_j:\Delta_j\to\mathbb C^2$:
$\Delta_j\times\mathbb C^2\to J'_j$ is a covering map and the disc is
simply connected. Local inverses of this covering show that the lift
is holomorphic. The deck-equivariance identity holds modulo an
actual lattice vector, so, after choosing this lift, it is
\[
 B_j(g_jz)-R_{g_j}(z)B_j(z)
  =\Pi_c(g_jz)\left(
       \frac{(1-a)v_j+n\epsilon_j\widehat\delta}{m_j}
          +\ell_j\right),\qquad \ell_j\in\Lambda.
 \tag{PS10}\label{pshift:eq:finite-equation}
\]
The integer column $\ell_j$ is constant: it is locally constant by
discreteness of the marked lattice and the disc is connected.
Iterating $m_j$ times, with $A_jv_j=v_j$ and
$A_j\widehat\delta=\widehat\delta$, yields the necessary equality
\[
 N_j\ell_j=-(1-a)v_j-n\epsilon_j\widehat\delta.
 \tag{PS11}\label{pshift:eq:finite-norm-obstruction}
\]
Here the left side of the iterated equation is zero because the
base point and the original vector coordinate return after $m_j$
steps; real invertibility of $\Pi_c$ makes the lattice equality exact.

If $a=-1$, applying $\gamma$ to (PS11) gives
\[
                   m_j\gamma(\ell_j)=-2\epsilon_j.
\]
This is impossible already at $j=1$, since $3$ does not divide $2$.
Thus a global isomorphism cannot have negative linear part.

For $a=1$, the first row of (PS11) gives $\gamma(\ell_j)=0$.
Equation (PS9) then proves that the order-three filling imposes no
restriction on $n$, whereas the order-four filling imposes precisely
$n\in2\mathbb Z$. More intrinsically, the finite obstruction is the
class of $-n\epsilon_j\widehat\delta$ in
\[
 \Lambda^{A_j}/N_j\Lambda.
 \tag{PS12}\label{pshift:eq:finite-class}
\]
In the original fixed bases this quotient is
$\mathbb Z/3$ at $j=1$, with $\widehat\delta$ zero, and
$\mathbb Z/4\oplus\mathbb Z/2$ at $j=2$, with
$\widehat\delta$ generating the second summand. Indeed (PS9) gives
\[
 N_1\Lambda=3\mathbb Z v_1\oplus\mathbb Z\widehat\delta,
 \quad
 N_2\Lambda=4\mathbb Z\varepsilon'\oplus
                      2\mathbb Z\widehat\delta,
 \qquad \varepsilon'=(1,3,-3,0)^{\mathsf t}.
\]
This proves the exact parity obstruction rather than merely failure
of the untranslated identity.

\subsubsection{Local maps for every even shift}

Let $n\in2\mathbb Z$. Define the following original rational columns,
integral cocycles, and holomorphic translation vectors:
\[
\begin{aligned}
 r_1&=-\frac n3(\widehat u+\widehat w),&
 \ell_1&=-n\widehat w,&
 b_1(z)&=\Pi_c(z)r_1=-\frac n3(\tau(z)+1,\mu(z))^{\mathsf t},\\
 r_2&= \frac n4(\widehat u+\widehat w),&
 \ell_2&=\frac n2\widehat u,&
 b_2(z)&=\Pi_c(z)r_2=\frac n4(\tau(z)+1,\mu(z))^{\mathsf t}.
\end{aligned}
\tag{PS13}\label{pshift:eq:local-corrections}
\]
Direct multiplication of the complete matrices in (PS3) gives
\[
 (I-A_j)r_j=\frac{n\epsilon_j}{m_j}\widehat\delta+\ell_j.
 \tag{PS14}\label{pshift:eq:local-matrix-identity}
\]
For $j=1$ the four entries are $(0,0,-n,n/3)$, and for $j=2$
they are $(0,n/2,0,-n/4)$. Pairing with the full $\Pi_c(g_jz)$
shows that $b_j$ satisfies (PS10) for $a=1$.
Translation by $b_j$ is therefore already an isomorphism of the
two finite local fillings. The remaining issue is its exact overlap
with the fixed middle identification.

Keep the original linearizing coordinate, radius, and phase
$s_j(g_jz)=e^{-2\pi i/m_j}s_j(z)$. On a logarithm branch put
\[
 w_j(z)=\frac{\log s_j(z)}{2\pi i},\qquad
 a_j(z)=b_j(z)+n\epsilon_j w_j(z)e.
 \tag{PS15}\label{pshift:eq:finite-middle-sections}
\]
Changing the logarithm adds $n\epsilon_j\widehat\delta$ under the
period map, so $a_j$ is branch independent as a torus section.
For the specified negative generator, $w_j(g_jz)=w_j(z)-1/m_j$.
Equations (PS4) and (PS14) consequently give the exact equality
\[
 a_j(g_jz)=R_{g_j}(z)a_j(z)+\Pi_c(g_jz)\ell_j.
 \tag{PS16}\label{pshift:eq:local-affine-cocycle}
\]
This is why the logarithmic term in (PS15) cannot be omitted.
In fact the original finite overlap maps are translations by
$\sigma_{j,c}=w_j\Pi_c v_j$, and
\[
                  \sigma_{j,c+n}-\sigma_{j,c}
                     =n\epsilon_jw_je.
 \tag{PS17}\label{pshift:eq:actual-overlap-difference}
\]
Thus $a_j$ is exactly the middle translation induced by the local
translation $b_j$, with the original gluing maps retained.

At the cusp retain $s=\tau-h$, $e^{2\pi is}=t_c$, $s(g_0z)=s(z)-1$,
and the complete period matrix
\[
 Z_c=sB_0+C_c(t_c),\quad
 B_0=\begin{pmatrix}0&1\\-1&0\end{pmatrix},\quad
 C_c=\begin{pmatrix}6\mu&h\\ b^{\mathrm{part}}-h+c&\mu\end{pmatrix}.
 \tag{PS18}\label{pshift:eq:cusp-periods}
\]
For $\lambda=(\lambda_1,\lambda_2)^{\mathsf t}\in\mathbb Z^2$,
$C_{c+n}\lambda-C_c\lambda=(0,n\lambda_1)^{\mathsf t}$. Therefore
\[
 \exp(2\pi iC_{c+n}\lambda)t_c^{B_0\lambda}
       =\exp(2\pi iC_c\lambda)t_c^{B_0\lambda}.
 \tag{PS19}\label{pshift:eq:cusp-actions-identical}
\]
The entire cusp actions, including their nonconstant holomorphic units,
are identical for an integer shift.

The affine cocycle convention in (PS16) has multiplication law
\[
                    \ell_{gh}=\ell_g+A_g\ell_h.
 \tag{PS20}\label{pshift:eq:cocycle-law}
\]
Since $g_0=(g_1g_2)^{-1}$, it gives
\[
\begin{aligned}
 \ell_1+A_1\ell_2
   &=-\frac{3n}{2}\widehat w+\frac n2\widehat\delta,\\
 \ell_0=-A_{g_0}(\ell_1+A_1\ell_2)
   &=\frac{3n}{2}\widehat w-\frac n2\widehat\delta.
\end{aligned}
\tag{PS21}\label{pshift:eq:cusp-cocycle}
\]
Here $A_{g_0}$ fixes both toric lattice vectors, as is also immediate
from $A_{g_0}=\left(\begin{smallmatrix}I&0\\B_0&I\end{smallmatrix}\right)$.
Put
\[
 q_n=(-3n/2,n/2)^{\mathsf t}\in\mathbb Z^2,
                  \qquad a_0(z)=s(z)q_n.
 \tag{PS22}\label{pshift:eq:cusp-translation}
\]
Then $a_0(g_0z)-a_0(z)=-q_n=\Pi_c(g_0z)\ell_0$, exactly as
required by (PS16). Exponentiating gives the map
$(x,t_c)\mapsto(t_c^{q_n}x,t_c)$ on the dense cusp torus.
It extends holomorphically and invertibly to the original infinite fan:
the integral map
\[
 (y,y_3)\longmapsto(y+y_3q_n,y_3)
\]
translates each height-one triangle by the integer vector $q_n$ and
therefore permutes the actual cones, with inverse given by $-q_n$.
It commutes with the cusp quotient action in (PS19). This proves
extension across every cusp stratum, not only across its dense torus.

\subsubsection{The global additive correction and all overlap identities}

We give the global correction explicitly as an additive gluing
construction and prove its solvability. Let $\widetilde B^\circ$ be
the universal cover of the three-punctured base. Its fundamental group
is free on the original two negative peripheral generators. Pull back
the original $\Pi_c,R$ to this cover. Assign $\ell_1,\ell_2$ from
(PS13) to its free generators and extend by (PS20). Freeness makes
this an actual integral cocycle; its third peripheral value is
precisely (PS21). The maps
\[
 (z,y)\longmapsto
       (gz,R_g(z)y+\Pi_c(gz)\ell_g)
 \tag{PS23}\label{pshift:eq:affine-bundle-action}
\]
form an affine action. To check the action law, covariance gives
$R_g(hz)\Pi_c(hz)=\Pi_c(ghz)A_g$; substituting this equality
in the composition gives exactly (PS20).
The quotient is an affine holomorphic bundle over $B^\circ$, since
the deck action on the base is free and properly discontinuous and
one may use evenly covered base opens as affine charts.

On a finite punctured disc, (PS16) says that $a_j$ is a section of
this affine bundle. On the cusp punctured disc, (PS21)--(PS22) say
the same for $a_0$. Define a sheaf $\mathscr H_n$ on $B$ as follows.
Away from the special points it is the sheaf of sections of (PS23).
Near $p_j$, a section is required to have the form $a_j+\eta_j$,
where $\eta_j:\Delta_j\to\mathbb C^2$ is holomorphic and obeys
$\eta_j(g_jz)=R_{g_j}(z)\eta_j(z)$. Near the cusp it is required
to have the form $a_0+\eta_0$, with $\eta_0(t_c)$ holomorphic on
the full cusp disc. These are compatible definitions on the punctured
overlaps by the equal affine cocycles just proved. They also prove
local nonemptiness, including at all three special points.

The differences of these local sections form a vector sheaf
$\mathscr V$. Its ordinary sections have covariance $\eta(gz)=R_g\eta(z)$;
at the two finite points and the cusp its extensions are precisely
the holomorphic extensions specified above. Projection to the first
coordinate gives an exact sequence on the original compact base
\[
 0\longrightarrow\mathcal O_B e\longrightarrow\mathscr V
        \longrightarrow\mathcal L\longrightarrow0,
                  \qquad \mathcal L\simeq\mathcal O_B(-p_0).
 \tag{PS24}\label{pshift:eq:vector-sheaf}
\]
Here $\mathcal L$ is exactly the homogeneous period sheaf from
\cref{prop:period-input-audit}, not a newly chosen extension. Indeed
the first diagonal entries in (PS4) are its factors $-1/\tau$ and
$1/\tau$, and at the cusp both extensions require ordinary
holomorphicity. The original function
\[
 \Theta=\frac{(E_4\circ\tau)^2(E_6\circ\tau)^{1/2}}
                 {\Delta_{\mathrm{mod}}\circ\tau}
\]
has those exact factors, orders two and one in $s_1,s_2$, and a
simple pole in $t_c$. Dividing by $\Theta$ identifies the ordinary
and finite stalks with $\mathcal O_B$, and the cusp stalk with
$\mathfrak m_{p_0}$. This recovers the displayed $\mathcal L$
with its original cusp condition. These transformation and order
calculations are the ones used to construct the retained $\mu$;
no topological conclusion about $X_c$ is involved.

For completeness, exactness on the finite stalks in (PS24) retains
the off-diagonal entry of $R_g$. Given an equivariant first component
$u(z)$, begin with the vector $(u(z),0)^{\mathsf t}$ and average it
under the finite group action on vector-valued holomorphic functions.
The resulting vector is equivariant, and its first component is
still $u$. Thus projection is surjective there. Its kernel is
exactly the invariant second-component functions, which are
holomorphic functions of $s_j^{m_j}=t_j$, multiplied by $e$.
The ordinary and cusp stalks are immediate. This also proves that
$\mathscr V$ is locally free of rank two: lift a local frame of
$\mathcal L$ and join it to $e$.

Both $H^1(B,\mathcal O_B)$ and
$H^1(B,\mathcal O_B(-p_0))$ vanish. One can compute these two
groups without a degree approximation: on the two standard affine
charts of $\mathbb P^1$, split a holomorphic Laurent series on
their overlap into its nonnegative and negative powers. For
$\mathcal O_B$ these are the two coboundary contributions, and
for $\mathcal O_B(-p_0)$ its transition factor shifts the second
contribution by one, still leaving no missing exponent. The positive
part converges on the full first affine chart and the transformed
negative part on the full second affine chart. This proves the two
vanishings, equivalently the usual additive Cousin calculation on
these charts. The exact sequence (PS24) now gives
\[
                             H^1(B,\mathscr V)=0.
 \tag{PS25}\label{pshift:eq:no-analytic-obstruction}
\]

Here is how to choose the correction entering the maps. Choose
local sections $h_\alpha$ of $\mathscr H_n$, taking the displayed
$a_j,a_0$ on sufficiently small special discs and affine charts
elsewhere. Their differences $d_{\alpha\beta}=h_\beta-h_\alpha$
are a Cech $1$-cocycle with values in $\mathscr V$. Project it to
$\mathcal L$ and solve that additive cocycle, using its just proved
$H^1$ vanishing. Lift the resulting local $\mathcal L$ sections to
$\mathscr V$ by the local lifting construction in (PS24). Subtracting
their coboundary from $d$ leaves a scalar $\mathcal O_Be$ cocycle.
Solve that scalar additive cocycle by its $H^1$ vanishing. This
produces sections $v_\alpha$ with
$d_{\alpha\beta}=v_\beta-v_\alpha$. The local sections
$h_\alpha-v_\alpha$ agree and define a global section
$a\in\mathscr H_n(B)$. This is a two-stage construction of the
required correction, not an unproved additional gluing condition.
It gives holomorphic vectors $\eta_j,\eta_0$ with
\[
 a=a_j+\eta_j\quad(j=1,2),\qquad a=a_0+\eta_0
                              \quad\hbox{at the cusp}.
 \tag{PS26}\label{pshift:eq:global-correction}
\]
Any choices made in the additive solutions give maps with the same
asserted properties. Their differences as sections of this affine
bundle lie in $H^0(B,\mathscr V)=\mathbb C e$: this follows from
(PS24), $H^0(B,\mathcal L)=0$, and $H^0(B,\mathcal O_B)=\mathbb C$.

Projecting $a$ modulo $\Pi_c\Lambda$ gives a holomorphic section of
the middle torus family because every $\ell_g$ in (PS23) is integral.
The promised three kinds of local maps are therefore
\[
\begin{aligned}
 F_{\mathrm{mid}}[z,\zeta]&=[z,\zeta+a(z)],\\
 F_j[z,\zeta]&=[z,\zeta+b_j(z)+\eta_j(z)]\quad(j=1,2),\\
 F_0[x,t_c]&=[\exp(2\pi i\eta_0(t_c))t_c^{q_n}x,t_c].
\end{aligned}
\tag{PS27}\label{pshift:eq:full-isomorphism}
\]
The first descends by (PS23) and equality of the integer-shift
lattices. The finite maps descend by (PS10), (PS14), and the exact
covariance of $\eta_j$. The cusp map descends by (PS19); its
holomorphic-unit factor extends by torus multiplication on every
original chart, and its integer factor extends by the fan map in
(PS22). Each map is a biholomorphism, with inverse given by the
negative translation or the inverse torus multiplier and fan map.

On a finite overlap, write the original gluing as
$h_{j,c}(\zeta)=\zeta+\sigma_{j,c}$. Then (PS17), (PS26), and
(PS27) give the equality of actual torus maps
\[
                     h_{j,c+n}F_j=F_{\mathrm{mid}}h_{j,c}.
 \tag{PS28}\label{pshift:eq:finite-overlap-identity}
\]
Both sides add
$\sigma_{j,c}+a=\sigma_{j,c+n}+b_j+\eta_j$.
On the cusp overlap the original map is exponentiation of $\zeta$;
(PS22) and (PS26) give
$\exp(2\pi ia)=t_c^{q_n}\exp(2\pi i\eta_0)$, exactly the
third line of (PS27). Thus that overlap also agrees, including
all nonconstant period units. The maps and their inverses glue to
a global biholomorphism $X_c\to X_{c+n}$ over $B$.

The choices can be made coherently for every parameter and every even
shift. Both initial cocycle columns in (PS13) have $\gamma=0$.
Because every $A_g$ preserves $\gamma$, (PS20) shows that
$\gamma(\ell_g)=0$ for every word. Therefore the affine term
$\Pi_c(gz)\ell_g$ in (PS23) is independent of $c$. The displayed
$a_j,a_0$, and the vector sheaf $\mathscr V$ are also independent
of $c$. Perform the additive construction once for $n=2$, keeping
one resulting section $a^{(2)}$. For $n=2k$ choose
$a^{(2k)}=k a^{(2)}$ and multiply its local correcting vectors by
$k$. All their affine cocycles and local formulas are then exactly
those for $n=2k$. The resulting maps satisfy the full group law
\[
 F_{2h,c+2k}\circ F_{2k,c}=F_{2(h+k),c},\qquad
                 F_{0,c}=\mathrm{id}_{X_c}.
 \tag{PS28a}\label{pshift:eq:coherent-shift-action}
\]
This follows by adding the actual translation vectors in the first
two lines of (PS27), and by multiplying their units and adding the
integer exponents in its third line. The formulas are jointly
holomorphic in $c$, since their corrections are independent of it.

\subsubsection{Intrinsic necessity and the resulting parameter space}

The analytically computed complete linear system
$|\omega_{X_c}^{-2}|$ is the original map $f_c$, by
\cref{thm:canonical-ring-recovery}. A biholomorphism preserves its
canonical bundle, the complete space of sections, and its evaluation
map. Consequently it induces $\alpha\in\operatorname{PGL}_2(\mathbb C)$
with $f_{c'}F=\alpha f_c$. The three special fibres have distinct
intrinsic types: the cusp is reduced and nonnormal, while the other
two have multiplicities three and four. All remaining fibres are
smooth reduced tori. Thus $\alpha$ fixes $p_0,p_1,p_2$ individually.
In the original coordinate, a fractional-linear map fixing
$\infty,0,1$ has zero lower-left and upper-right entries and equal
diagonal entries; hence $\alpha=\mathrm{id}_B$.

Now \cref{pshift:lem:centralizer} forces $c'-c=n\in\mathbb Z$
and linear part $\pm I_2$. Equations (PS10)--(PS11), obtained on
the original normalized finite covers, exclude $-I_2$ and then
exclude odd $n$ for $+I_2$. Conversely (PS27)--(PS28) construct
the global map for every even $n$. This proves
\cref{pshift:thm:classification} without converting a failed
particular map into a nonisomorphism assertion.

As a complex parameter space for the isomorphism classes represented
by this retained family, the quotient $\mathcal P/(2\mathbb Z)$ has
the explicit coordinate
\[
 w=\exp(-\pi ic),\qquad 0<|w|<\exp(-\pi M).
 \tag{PS29}\label{pshift:eq:parameter-disc}
\]
Indeed $|w|=\exp(\pi\operatorname{Im}c)$, and two complex numbers
have the same value of this exponential exactly when they differ
by an even integer. Every point in the displayed punctured disc
has a logarithm giving a $c\in\mathcal P$. The exponential has
nonzero derivative, so these local inverse logarithms give the
claimed quotient complex structure. The integer-period quotient
seen by the middle family instead has coordinate $w^2$; it is
covered two-to-one by (PS29). Its two lifts are distinguished by
the exact order-two summand in (PS12). This statement classifies
the parameters of the retained family; it does not assert that
the family exhausts all deformations or all complex threefolds
with the same singular-fibre types.

